\documentclass[12pt,a4paper]{article}
\usepackage[utf8]{inputenc}
\usepackage[spanish]{babel}
\usepackage{geometry}
\usepackage{amsmath}
\usepackage{amssymb}
\usepackage{booktabs}
\usepackage{longtable}
\usepackage{graphicx}
\usepackage{enumitem}
\usepackage{float}
\usepackage{xcolor}
\usepackage{fancyhdr}
\usepackage{caption}
\usepackage[numbers]{natbib}
\usepackage[colorlinks=true,linkcolor=black,citecolor=blue,urlcolor=blue]{hyperref}

\geometry{margin=2.5cm}

\pagestyle{fancy}
\fancyhf{}
\fancyhead[L]{Sistema de Tratamiento ABR/RAP, TAF,SS}
\fancyhead[R]{\thepage}

\setcounter{tocdepth}{2}

\begin{document}

\begin{titlepage}
\centering
\vspace*{2cm}

{\Huge\bfseries Memoria de Cálculo\par}
\vspace{1.5cm}

{\Large\itshape Dimensionamiento de:\par}
\vspace{0.5cm}

{\LARGE\bfseries Sistema de Tratamiento ABR/RAP, TAF,SS\par}
\vspace{2cm}

{\large Caudal de diseño: 17.31 L/s\par}
{\large Número de líneas: 3\par}
\vfill
{\large\today\par}
\end{titlepage}

\tableofcontents

\listoffigures

\listoftables

\newpage
\section{Resumen del Proyecto}

Este documento presenta el dimensionamiento detallado del Sistema de Tratamiento ABR/RAP, TAF,SS, disenado para tratar un caudal de 17.31 L/s distribuido en 3 lineas paralelas (5.77 L/s por linea).

\textbf{Caracteristicas del afluente:}
\begin{itemize}
    \item DBO$_5$: 243.1 mg/L
    \item DQO: 498 mg/L
    \item SST: 156.0 mg/L
    \item Coliformes fecales: 3e+03 NMP/100mL
    \item Temperatura: 25.6°C
\end{itemize}

\textbf{Secuencia de unidades del tren:} 
Canal de Desbaste con Rejillas $\rightarrow$ Desarenador de Flujo Horizontal $\rightarrow$ Reactor Anaerobio con Pantallas (ABR/RAP) $\rightarrow$ Biofiltro de Carga Mecanizada e Hidráulica $\rightarrow$ Sedimentador Secundario $\rightarrow$ Sistema de Desinfeccion con Hipoclorito de Sodio $\rightarrow$ Lechos de Secado de Lodos

\newpage
\section{Dimensionamiento de Unidades}

\subsection{Canal de Desbaste con Rejillas}

\

Las rejillas constituyen la primera barrera de proteccion del sistema, reteniendo solidos gruesos como plasticos, ramas y papel que podrian danar equipos o causar obstrucciones en tuberias aguas abajo. El diseno hidraulico de esta unidad debe garantizar velocidades suficientes para arrastrar solidos sedimentables, pero no tan elevadas que dificulten el paso del agua a traves de las barras.

Segun los criterios establecidos por Metcalf y Eddy \cite{metcalf2014}, las velocidades de diseno en canales con rejillas deben mantenerse entre 0.40 y 0.60~m/s. Para este proyecto se adopta un valor intermedio de 0.50~m/s, lo cual resulta apropiado considerando el caudal de diseno. El tirante hidraulico se establece en 0.40~m, valor que permite una velocidad de flujo uniforme en la seccion del canal y evita la sedimentacion de solidos organicos antes de la rejilla, segun los criterios de Metcalf y Eddy \cite{metcalf2014}.

La perdida de carga en rejillas limpias se calcula mediante la formula de Kirschmer (1926), que relaciona la geometria de las barras con las caracteristicas del flujo. Los parametros que determinan esta perdida son: el espaciado entre barras ($b$), el espesor de las barras ($w$), la velocidad del flujo ($v$) y el angulo de inclinacion ($\theta$):

\begin{equation}
h_L = \beta \cdot \left(\frac{w}{b}\right)^{4/3} \cdot \frac{v^2}{2g} \cdot \sin\theta
\end{equation}
\captionequation{Perdida de carga en rejillas limpias - Formula de Kirschmer}

\textit{Donde:}
\begin{itemize}[noitemsep,leftmargin=2em]
    \item[$h_L$] = Perdida de carga en rejillas limpias (m)
    \item[$\beta$] = Factor de forma de la barra (2.42 para barras rectangulares)
    \item[$w$] = Espesor de la barra (10 mm)
    \item[$b$] = Espaciado entre barras (15 mm)
    \item[$v$] = Velocidad del flujo en el canal (m/s)
    \item[$g$] = Aceleracion de la gravedad (9,81 m/s\textsuperscript{2})
    \item[$\theta$] = Angulo de inclinacion de las barras (60\textdegree)
\end{itemize}

\begin{table}[H]
\centering
\caption{Parametros de diseno -- Rejillas}
\begin{tabular}{llcc}
\toprule
Parametro & Descripcion & Valor adoptado & Fuente \\
\midrule
Velocidad en canal & Velocidad de paso & 0.50 m/s & Metcalf y Eddy \\
Tirante hidraulico & Altura de agua en canal & 0.40 m & Criterio tecnico \\
Angulo de rejilla & Inclinacion barras & 60\textdegree & Practica comun \\
Espaciado entre barras & Abertura libre & 15 mm & OPS/CEPIS \\
Espesor de barras & Seccion barra & 10 mm & Constructivo \\
Factor de forma & Barras rectangulares & 2.42 & Kirschmer \\
\bottomrule
\end{tabular}
\end{table}

El caudal de diseno por linea es 5.8 L/s. La seccion transversal del canal se determina aplicando la ecuacion de continuidad:

\begin{equation}
A_{{canal}} = \frac{Q}{v} = \frac{0.00577}{0.50} = 0.01154 \text{ m}^2
\end{equation}

\textit{Donde:}
\begin{itemize}[noitemsep,leftmargin=2em]
    \item[$A_{{canal}}$] = Seccion transversal del canal (m\textsuperscript{2})
    \item[$Q$] = Caudal de diseno por linea (0.00577 m\textsuperscript{3}/s)
    \item[$v$] = Velocidad de diseno adoptada (0.50 m/s)
\end{itemize}

Despejando el ancho del canal de la ecuacion $A = b \times h$, con $h = 0.40$ m:

\begin{equation}
b = \frac{A_{{canal}}}{h} = \frac{0.01154}{0.40} = 0.029 \text{ m}
\end{equation}

\textit{Donde:}
\begin{itemize}[noitemsep,leftmargin=2em]
    \item[$b$] = Ancho teorico del canal (m)
    \item[$A_{{canal}}$] = Seccion transversal del canal (0.01154 m\textsuperscript{2})
    \item[$h$] = Tirante hidraulico (0.40 m)
\end{itemize}

Este ancho teorico de 0.029 m es inferior al minimo constructivo practico. Se adopta un ancho de 0.60 m como dimension minima constructiva segun criterios de OPS/CEPIS \cite{ops2005} y SENAGUA \cite{senagua2012}, que establecen un ancho minimo de 0.60 m para permitir la operacion adecuada, el acceso para limpieza y la aplicacion de herramientas de mantenimiento. Verificando la velocidad real con este ancho:

\begin{equation}
v_{{real}} = \frac{Q}{A_{{real}}} = \frac{0.00577}{0.60 \times 0.40} = 0.024 \text{ m/s}
\end{equation}

\textit{Donde:}
\begin{itemize}[noitemsep,leftmargin=2em]
    \item[$v_{{real}}$] = Velocidad real verificada con ancho constructivo (m/s)
    \item[$Q$] = Caudal de diseno (0.00577 m\textsuperscript{3}/s)
    \item[$A_{{real}}$] = Seccion real con ancho constructivo adoptado (0.60 $\times$ 0.40 m\textsuperscript{2})
\end{itemize}

La velocidad real de 0.024 m/s es el resultado de adoptar el ancho minimo constructivo. La perdida de carga calculada con esta velocidad real resulta $h_L = 0.0036$ cm, menor al umbral de 15 cm que Metcalf y Eddy \cite{metcalf2014} indican como referencia para sistemas de limpieza mecanica.

La longitud de 2.0 m y el ancho de 0.60 m responden a criterios constructivos y operativos. El ancho de 0.60 m es la dimension minima constructiva segun OPS/CEPIS \cite{ops2005} y SENAGUA \cite{senagua2012}, que permite la operacion adecuada de la rejilla, el acceso para limpieza y la aplicacion de herramientas de mantenimiento. La longitud contempla el espacio de la rejilla propiamente dicha mas las zonas de transicion necesarias para la operacion.





\subsubsection{Verificacion}

\textbf{Verificacion para Caudal Maximo Horario}

Se verifica el comportamiento hidraulico para el caudal maximo horario, aplicando un factor de pico tipico de 2.8 sobre el caudal medio:

\begin{equation}
Q_{{max}} = 2.8 \times Q_{{medio}} = 2.8 \times 5.8 = 16.0 \text{ L/s}
\end{equation}

A caudal maximo, la velocidad horizontal resulta:

\begin{equation}
v_{{max}} = \frac{Q_{{max}}}{A_{{real}}} = \frac{15.983/1000}{0.60 \times 0.40} = 0.067 \text{ m/s}
\end{equation}

\textbf{Criterios de verificacion de velocidad} (segun Metcalf y Eddy \cite{metcalf2014}):

La velocidad calculada $v_{{max}} = 0.067$~m/s se evalua segun:

\begin{equation}
\text{Estado} = \begin{cases}
    \text{OPTIMO} & \text{si } v_{{max}} \leq 1.5 \text{ m/s} \quad \text{(sin riesgo)} \\
    \text{ACEPTABLE} & \text{si } 1.5 < v_{{max}} \leq 2.0 \text{ m/s} \quad \text{(monitoreo)} \\
    \text{NO ADMISIBLE} & \text{si } v_{{max}} > 2.0 \text{ m/s} \quad \text{(dano seguro)}
\end{cases}
\end{equation}
\captionequation{Criterios de verificacion de velocidad en rejillas}

Con la velocidad calculada $v_{{max}} = 0.067$~m/s, la unidad se clasifica como \textbf{OPTIMO}. En terminos de cumplimiento, la verificacion de velocidad \textbf{CUMPLE} porque es menor que el limite recomendado de 1.5~m/s para operacion segura.

Respecto a la perdida de carga, el valor maximo calculado es 0.0276~cm. Por tanto, la verificacion de perdida de carga \textbf{CUMPLE}, ya que es menor que el umbral de 15~cm para limpieza manual aceptable.

\subsubsection{Resultados}

\begingroup
\small
\begin{longtable}{lccc}
\caption{Verificacion de criterios de diseno - Rejillas}\\
\toprule
Parametro & Valor calculado & Criterio & Estado \\
\midrule
\endfirsthead
\caption[]{Verificacion de criterios de diseno - Rejillas (continuación)}\\
\toprule
Parametro & Valor calculado & Criterio & Estado \\
\midrule
\endhead
\midrule
\multicolumn{4}{r}{\textit{Continúa en la siguiente página}} \\
\endfoot
\bottomrule
\endlastfoot
Velocidad de diseno & 0.50 m/s & 0.40 -- 0.60 m/s & Cumple \\
Velocidad real & 0.024 m/s & Ancho minimo constructivo$^a$ & Aplicado \\
Perdida de carga (Qmax) & 0.0276 cm & $<$ 15 cm & Cumple \\
Ancho constructivo & 0.60 m & $\geq$ 0.30 m & Cumple \\
\multicolumn{4}{p{\textwidth}}{\small $^a$ Para caudales pequenos (5.8 L/s) el ancho minimo constructivo (0.60 m) domina sobre el criterio de velocidad.} \\
\end{longtable}
\endgroup

\begin{figure}[H]
\centering
\includegraphics[width=0.9\textwidth]{figuras/rejillas_rejillas.png}
\caption{Esquema del canal de desbaste con rejillas (dimensiones en metros)}
\label{fig:rejillas}
\end{figure}





\textit{Nota sobre cantidades:} El presente dimensionamiento no estima la masa o volumen diario de residuos de cribado. Segun Metcalf y Eddy \cite{metcalf2014}, la produccion tipica de residuos de rejillas finas en aguas residuales municipales varia entre 0.005--0.02 m\textsuperscript{3}/1000 m\textsuperscript{3} de agua tratada. Para el caudal de diseno, esto representaria aproximadamente 0.02--0.09 m\textsuperscript{3}/d (total planta), sujeto a caracterizacion local.

Las dimensiones finales del canal con rejillas son: ancho 0.60 m, tirante 0.40 m y longitud 2.0 m. Se requieren dos unidades, una por cada linea de tratamiento.

\subsection{Desarenador de Flujo Horizontal}



El desarenador remueve part\'iculas minerales (arena, grava) con di\'ametro superior a 0.20~mm mediante sedimentaci\'on gravitacional. Seg\'un Metcalf y Eddy \cite{metcalf2014}, los par\'ametros de dise\~no son: velocidad horizontal 0.25--0.30~m/s (m\'aximo 0.30~m/s para evitar arrastre de sedimentos), tiempo de retenci\'on 30--60~s, y profundidad de flujo 0.75--2.0~m (OPS/CEPIS \cite{ops2005} permite valores menores en plantas peque\~nas).

La profundidad total del canal incluye: profundidad \'util de flujo (0.40~m), zona de almacenamiento de arena (0.25--0.30~m; se adopta 0.30~m), y bordo libre (0.30~m).

El dise\~no del desarenador se fundamenta en la teor\'ia de sedimentaci\'on discreta. Primero se calcula la velocidad de sedimentaci\'on de las part\'iculas objetivo mediante la Ley de Stokes:

\begin{equation}
v_s = \frac{g \cdot (S_s - 1) \cdot d^2}{18 \cdot \nu}
\end{equation}
\captionequation{Ley de Stokes - Velocidad de sedimentacion de particulas}

\textit{Donde:}
\begin{itemize}[noitemsep,leftmargin=2em]
    \item[$v_s$] = Velocidad de sedimentaci\'on (m/s)
    \item[$g$] = Aceleraci\'on de la gravedad (9.81 m/s²)
    \item[$S_s$] = Gravedad espec\'ifica de la arena (2.65)
    \item[$d$] = Di\'ametro de la part\'icula (0.00020 m)
    \item[$\nu$] = Viscosidad cinem\'atica del agua a 25.6°C ($8.3973e-07$ m²/s)
\end{itemize}

\begin{equation}
v_s = \frac{9.81 \times (2.65 - 1) \times (0.00020)^2}{18 \times 8.3973e-07} = 0.043 \text{ m/s}
\end{equation}

Con este valor de $v_s$ se determina la longitud te\'orica del desarenador mediante la ecuaci\'on de sedimentaci\'on tipo I:

\begin{equation}
L = \frac{H \cdot v_h}{v_s} = \frac{0.40 \times 0.30}{0.043} = 2.8 \text{ m}
\end{equation}

\textit{Donde:}
\begin{itemize}[noitemsep,leftmargin=2em]
    \item[$L$] = Longitud del desarenador (m)
    \item[$H$] = Profundidad \'util (0.40 m)
    \item[$v_h$] = Velocidad horizontal adoptada (0.30 m/s)
\end{itemize}

La longitud te\'orica resultante (2.8~m) es adecuada para el caudal de dise\~no y se adopta como longitud de dise\~no.

El ancho del canal se determina por la ecuaci\'on de continuidad:

\begin{equation}
B = \frac{Q}{v_h \cdot H} = \frac{0.00577}{0.30 \times 0.40} = 0.048 \text{ m}
\end{equation}

\textit{Donde:}
\begin{itemize}[noitemsep,leftmargin=2em]
    \item[$B$] = Ancho te\'orico del canal (m)
    \item[$Q$] = Caudal de dise\~no (0.00577 m³/s)
\end{itemize}

El ancho te\'orico resulta inferior al m\'inimo constructivo. Se adopta 0.60 m seg\'un OPS/CEPIS \cite{ops2005} y SENAGUA \cite{senagua2012}, que establecen 0.60 m como ancho m\'inimo para operaci\'on y mantenimiento.

Verificando la velocidad horizontal real con el ancho adoptado:

\begin{equation}
v_{real} = \frac{Q}{B \cdot H} = \frac{0.00577}{0.60 \times 0.40} = 0.024 \text{ m/s}
\end{equation}

El tiempo de retenci\'on real resulta:

\begin{equation}
t_r = \frac{L}{v_{real}} = \frac{3.0}{0.024} = 124.8 \text{ s}
\end{equation}

\subsubsection{Verificacion}

\textbf{Verificacion de Velocidad Critica de Resuspension}

Para garantizar que la arena sedimentada no sea resuspendida por el flujo, se verifica que la velocidad horizontal del dise\~no sea menor que la velocidad cr\'itica de resuspensi\'on, calculada mediante la f\'ormula de Camp-Shields:

\begin{equation}
v_c = \sqrt{\frac{8 \beta g (S_s - 1) d}{f}}
\end{equation}
\captionequation{Velocidad critica de resuspension - Camp-Shields}

\textit{Donde:}
\begin{itemize}[noitemsep,leftmargin=2em]
    \item[$v_c$] = Velocidad cr\'itica de resuspensi\'on (m/s)
    \item[$\beta$] = Factor de forma (0.05, rango 0.04--0.06)
    \item[$g$] = Aceleraci\'on de la gravedad (9.81 m/s²)
    \item[$S_s$] = Gravedad espec\'ifica de la arena (2.65)
    \item[$d$] = Di\'ametro de la part\'icula (0.00020 m)
    \item[$f$] = Factor de fricci\'on de Darcy-Weisbach (0.025, rango 0.02--0.03)
\end{itemize}

\begin{equation}
v_c = \sqrt{\frac{8 \times 0.05 \times 9.81 \times (2.65 - 1) \times 0.00020}{0.025}} = 0.228 \text{ m/s}
\end{equation}
\captionequation{Velocidad critica calculada}

\textbf{Verificacion para Caudal Maximo Horario}

Se verifica el comportamiento hidr\'aulico para el caudal m\'aximo horario, aplicando un factor de pico t\'ipico de 2.8 sobre el caudal medio:

\begin{equation}
Q_{max} = 2.8 \times Q_{medio} = 2.8 \times 5.8 = 16.0 \text{ L/s}
\end{equation}

A caudal m\'aximo, la velocidad horizontal resulta:

\begin{equation}
v_{h,max} = \frac{Q_{max}}{B \cdot H} = \frac{15.983/1000}{0.60 \times 0.40} = 0.067 \text{ m/s}
\end{equation}

El tiempo de retenci\'on a caudal m\'aximo:

\begin{equation}
t_{r,max} = \frac{L}{v_{h,max}} = \frac{3.0}{0.067} = 45.0 \text{ s}
\end{equation}

\textbf{Criterios de verificaci\'on de resuspensi\'on} (Camp-Shields, Metcalf y Eddy \cite{metcalf2014}):

La velocidad calculada $v_{h,max} = 0.067$~m/s se eval\'ua seg\'un (l\'imite: $v_c \times 90\% = 0.205$~m/s):

\begin{equation}
\text{Estado} = \begin{cases}
    \text{SEGURIDAD} & \text{si } v_{h,max} \leq 0.205 \text{ m/s} \quad \text{(no hay resuspensi\'on)} \\
    \text{RIESGO} & \text{si } v_{h,max} > 0.205 \text{ m/s} \quad \text{(arena resuspendida)}
\end{cases}
\end{equation}
\captionequation{Criterios de verificacion de resuspension}

Con la velocidad calculada $v_{h,max} = 0.067$~m/s, la unidad se clasifica como \textbf{SEGURIDAD}. En terminos de cumplimiento, la verificacion de resuspension \textbf{CUMPLE} porque es menor que el limite de resuspension con factor de seguridad (0.205~m/s), por lo que no se produce resuspension de la arena depositada.

\subsubsection{Resultados}

\begingroup
\small
\begin{longtable}{l p{9cm}}
\caption{Dimensiones del desarenador}\\
\toprule
Par\'ametro & Valor \\
\midrule
\endfirsthead
\caption[]{Dimensiones del desarenador (continuación)}\\
\toprule
Par\'ametro & Valor \\
\midrule
\endhead
\midrule
\multicolumn{2}{r}{\textit{Continúa en la siguiente página}} \\
\endfoot
\bottomrule
\endlastfoot
Longitud dise\~no & 3.0 m (m\'inimo constructivo SENAGUA) \\
Ancho & 0.60 m (m\'inimo constructivo) \\
Profundidad \'util & 0.40 m \\
Velocidad horizontal real & 0.024 m/s \\
Velocidad cr\'itica resuspensi\'on & 0.228 m/s (Camp-Shields) \\
Tiempo retenci\'on real & 124.8 s \\
\multicolumn{2}{p{\textwidth}}{\small\textit{$^a$Nota: El TRH de 125~s y velocidad de 0.024~m/s se deben al ancho m\'inimo constructivo de 0.60~m. }} \\
\end{longtable}
\endgroup

\textbf{Manejo de arena removida}

El desarenador retiene part\'iculas minerales, principalmente arena y grava fina, que se sedimentan en el fondo del canal por acci\'on gravitacional. Estas part\'iculas presentan di\'ametros superiores a 0.20~mm, con gravedad espec\'ifica de 2.65, correspondiendo a s\'olidos inorg\'anicos de alta densidad presentes en el agua residual. El canal incorpora una zona de almacenamiento de 0.30~m de altura, dimensionada espec\'ificamente para acumular la arena removida entre operaciones de limpieza. La frecuencia de limpieza recomendada es diaria o seg\'un la acumulaci\'on observada, manteniendo la arena acumulada por debajo del 50\% de la altura de almacenamiento designada para garantizar el funcionamiento \'optimo del sistema. La remoci\'on se realiza manualmente mediante palas y cubos desde la plataforma de operaci\'on, aunque puede implementarse un sistema de succi\'on o bombeo si el caudal y la producci\'on de arena lo justifican t\'ecnicamente.

\textbf{Disposici\'on final de la arena:}

La arena removida en el desarenador puede manejarse de diferentes formas seg\'un su composici\'on y la normativa local vigente. La opci\'on m\'as com\'un es el relleno sanitario, ya que la arena es material inorg\'anico inerte que puede disposicionarse en rellenos autorizados sin restricciones especiales. En caso de que la arena est\'e relativamente limpia, con bajo contenido de materia org\'anica adherida, existe la posibilidad de reuso condicionado para relleno y compactaci\'on en obras civiles previa caracterizaci\'on, mezcla con concreto para elementos no estructurales, o como complemento en sistemas de filtraci\'on gruesa previo lavado. Si se detecta alta carga org\'anica adherida a los granos de arena, evidenciada por olor f\'etido o color oscuro, se recomienda un lavado previo antes de la disposici\'on final para reducir el impacto ambiental.

\begin{figure}[H]
\centering
\includegraphics[width=\textwidth]{figuras/Esquema_Desarenador.png}
\caption{Esquema del desarenador de flujo horizontal con perfil longitudinal y corte transversal. Se muestran la entrada del afluente, la zona de flujo, el dep\'osito de arena y la salida . }
\label{fig:desarenador}
\end{figure}

\textit{Nota sobre cantidades:} El presente dimensionamiento no estima la masa o volumen diario de arena removida. Seg\'un Metcalf y Eddy \cite{metcalf2014}, la producci\'on t\'ipica de arena en aguas residuales municipales var\'ia entre 0.004--0.15 L/m\textsuperscript{3} de agua tratada, dependiendo de las caracter\'isticas del sistema de alcantarillado (red combinada vs. separada) y del suelo local. Para el caudal de dise\'no, esto representar\'ia aproximadamente 0.02--0.6 L/d (total planta), sujeto a caracterizaci\'on local.

\subsection{Reactor Anaerobio con Pantallas (ABR/RAP)}

\subsubsection{Dimensionamiento}

\textbf{Nomenclatura del reactor:} En la literatura técnica internacional, el reactor objeto de este manual se denomina ABR (\textit{Anaerobic Baffled Reactor}). En la literatura hispanohablante y en particular en los manuales de la OPS/OMS para saneamiento en países de desarrollo, se utiliza con frecuencia la sigla RAP (Reactor Anaerobio con Pantallas), que traduce literalmente el mismo concepto. Este documento adopta la denominación ABR como nombre principal del documento y del módulo de cálculo, por ser el término de uso universal en la bibliografía técnica de referencia. La sigla RAP se utiliza en paralelo en los títulos y encabezados del documento —de ahí la denominación compuesta "ABR / RAP"— con el fin de garantizar la trazabilidad con los manuales de la OPS/OMS que sirven de referencia complementaria en este proyecto.

\textbf{¿Qué es un Reactor Anaerobio con Pantallas (ABR / RAP)?}

El Reactor Anaerobio con Pantallas —ABR por sus siglas en inglés— es un sistema de tratamiento biológico anaerobio de aguas residuales que opera mediante flujo a través de una serie de compartimentos separados por pantallas o tabiques verticales (bafles). El agua residual entra por un extremo del reactor y recorre, en serie, cada uno de los compartimentos hasta llegar al extremo opuesto donde descarga el efluente tratado. En cada compartimento, el flujo sigue un patrón alternado de ascenso y descenso determinado por la posición y configuración de las pantallas.

El ABR fue desarrollado y descrito formalmente por McCarty y colaboradores en la Universidad de Stanford (Bachmann et al., 1985) como una evolución simplificada de los reactores anaerobios de lecho de lodos. Su gran ventaja conceptual sobre el reactor UASB (\textit{Upflow Anaerobic Sludge Blanket}) es la ausencia de un separador trifásico: la retención de la biomasa se logra por medios puramente hidráulicos y gravitacionales, aprovechando el hecho de que el flujo ascendente en cada compartimento es suficientemente lento como para mantener el manto de lodos en posición estable.

El ABR ocupa una posición intermedia entre el tanque séptico y el UASB en cuanto a complejidad constructiva y eficiencia de tratamiento. Es más eficiente que el tanque séptico porque el tratamiento ocurre en múltiples etapas en serie, cada una con su propia comunidad microbiana especializada. Es más simple que el UASB porque no requiere el separador trifásico, que es el componente más costoso y delicado de ese reactor. Esta combinación de buena eficiencia y simplicidad constructiva lo convierte en una opción especialmente adecuada para sistemas de saneamiento descentralizados, comunidades rurales y pequeñas cabeceras parroquiales en países de ingresos medios y bajos.

\textbf{Rol en el tren de tratamiento del proyecto:}

En el presente proyecto, el ABR se diseña como unidad de tratamiento primario anaerobio, situada aguas abajo del pretratamiento (cribado y desarenado) y aguas arriba de las unidades de tratamiento complementario o secundario. Su función principal es la remoción de materia orgánica en condiciones anaerobias, la estabilización parcial de los sólidos sedimentables y la acumulación y digestión \textit{in situ} de los lodos generados.

En la configuración típica del proyecto, el ABR reemplaza o complementa al tanque séptico convencional en situaciones donde se requiere un mayor nivel de remoción de DBO y DQO en la etapa anaerobia, con mayor control del proceso y sin la necesidad de un sistema mecánico de aireación. El efluente del ABR puede destinarse a un campo de infiltración, a un humedal de flujo subsuperficial, a un filtro anaerobio de flujo ascendente (FAFA) o, según el nivel de tratamiento requerido, a una etapa de tratamiento aerobio posterior (biofiltro, lodos activados, etc.).

\textbf{Principio de funcionamiento por compartimentos:}

El ABR es un reactor rectangular dividido internamente en una serie de compartimentos por medio de tabiques o pantallas verticales. Los tabiques no llegan completamente al fondo ni completamente a la superficie: dejan aberturas alternadas para forzar el recorrido del agua a través de toda la altura de cada compartimento.

En la configuración adoptada en este proyecto, los tabiques tienen una abertura en la parte inferior. El agua entra por el primer compartimento por el fondo, sube a través de él, desborda por encima del tabique intermedio o pasa por una abertura superior, desciende hasta el fondo del segundo compartimento, vuelve a subir, y así sucesivamente hasta el último compartimento. Este patrón de flujo ascendente-descendente garantiza que el agua atraviese toda la profundidad del manto de lodos en cada compartimento, maximizando el contacto entre el sustrato y la biomasa.

\textbf{Retención de biomasa sin separador trifásico:}

La clave del funcionamiento del ABR es la velocidad ascensional relativamente baja en el interior de cada compartimento. El manto de lodos anaerobio se mantiene en posición porque la velocidad ascendente del flujo no es suficiente para arrastrar los gránulos o flóculos de lodo hacia el compartimento siguiente. Esta retención es fundamentalmente diferente a la del UASB, donde se requiere un separador trifásico para capturar los sólidos que el flujo ascendente tiende a arrastrar hacia la superficie: en el ABR, la propia geometría del sistema asegura que la velocidad ascensional se mantenga por debajo del umbral de arrastre del lodo, sin necesidad de separadores mecánicos adicionales.

La consecuencia práctica es que en el ABR cada compartimento acumula su propia población microbiana, y esa población se adapta progresivamente a la calidad del agua que recibe: el primer compartimento, que recibe el agua cruda más concentrada, contiene predominantemente bacterias hidrolíticas y acidogénicas; los compartimentos intermedios tienen una mezcla de acidogénicos y acetogénicos; y el último compartimento, que recibe un agua ya considerablemente tratada, tiene condiciones más propicias para los metanogénicos. Esta separación funcional a lo largo del reactor es uno de los factores que explican la robustez del ABR frente a variaciones de carga.

\textbf{Mecanismos principales de remoción:}

Los mecanismos de remoción en el ABR son los propios del tratamiento anaerobio en cuatro etapas clásicas:

\textit{Hidrólisis:} Las moléculas orgánicas complejas (proteínas, grasas, hidratos de carbono) son fraccionadas en monómeros y dímeros por enzimas extracelulares producidas por las bacterias hidrolíticas. Esta etapa es generalmente la limitante del proceso para residuales con alta concentración de sólidos particulados.

\textit{Acidogénesis:} Los productos de la hidrólisis son fermentados por bacterias acidogénicas a ácidos grasos volátiles (AGV), alcoholes, dióxido de carbono e hidrógeno. Esta etapa es generalmente rápida en comparación con la hidrólisis y con la metanogénesis.

\textit{Acetogénesis:} Los ácidos grasos volátiles de cadena corta (propionato, butirato) son convertidos a acetato, CO$_2$ e H$_2$ por bacterias acetogénicas. Esta etapa es termodinámicamente desfavorable a concentraciones altas de H$_2$ y solo ocurre de forma eficiente cuando el hidrógeno es consumido rápidamente por organismos hidrogenotróficos (sintrofía obligada).

\textit{Metanogénesis:} El acetato y el H$_2$/CO$_2$ son convertidos a metano y CO$_2$ por las arqueas metanogénicas. Esta etapa es la más sensible a variaciones de pH y temperatura, y la que limita la tasa global del proceso cuando el sistema está correctamente operado.

Adicionalmente, en el ABR ocurre sedimentación de sólidos suspendidos por efecto de la reducción de velocidad dentro de cada compartimento: los sólidos gruesos que no se degradan rápidamente sedimentan en el fondo del primer compartimento o de los primeros compartimentos, donde se acumulan y se digieren lentamente. Esta función es análoga a la del tanque séptico, pero con la ventaja de que el tratamiento biológico activo ocurre en paralelo y de forma mucho más eficiente.

\textbf{Criterios de diseño adoptados por el proyecto:}

Los siguientes parámetros han sido seleccionados como criterios de diseño para el reactor ABR, basados en la bibliografía técnica y adaptados a las condiciones del proyecto:

\begin{table}[H]
\centering
\caption{Parámetros seleccionados para el diseño ABR / RAP}
\begin{tabular}{p{6cm}cc}
\toprule
\textbf{Parámetro} & \textbf{Valor adoptado} & \textbf{Rango recomendado} \\
\midrule
Modo de diseño & COV & TRH / COV / SRT \\
Tiempo de retención hidráulico (TRH) & 6 h & 24--72 h \\
TRH resultante & 5.5 h & $\geq$ 6.0 h \\
Número de compartimentos ($n$) & 4 & 2--8 (recomendado: 4--6) \\
Número de módulos paralelos & 3 & 1 o más \\
Ancho del módulo ($W$) & 1.79 m & $\leq$ 6.0 m \\
Ancho total de la línea & 6.36 m & --- \\
Relación $L_{comp}/W$ objetivo & 1.5 & --- \\
Velocidad ascensional de diseño ($v_{up}$) & 1.0 m/h & 0,2--1,5 m/h \\
Profundidad útil del líquido ($H$) & 2.0 m & 1,5--3,0 m \\
Zona de acumulación de lodos & 0.30 m & 0,2--0,5 m \\
Bordo libre & 0.30 m & $\geq$ 0,30 m \\
\bottomrule
\end{tabular}
\end{table}

\textit{Nota: Estos valores son criterios adoptados por el proyecto basados en Barber \& Stuckey (1999), Foxon et al. (2004) y OPS/OMS (2005). No son límites normativos impuestos por la TULSMA ni por normas ecuatorianas específicas para reactores anaerobios.}

\textbf{Ecuaciones de dimensionamiento:}

El dimensionamiento del ABR se fundamenta en el criterio de la carga orgánica volumétrica (COV), de modo que el reactor mantenga una tasa de carga compatible con el mantenimiento de una biomasa anaerobia estable. Las ecuaciones fundamentales son las siguientes:

\textbf{{[Ecuación 1] Volumen total del reactor (criterio COV):}}

\begin{equation}
V_{total} = \frac{Q_{medio} \times DBO_{entrada} / 1\,000}{COV_{dise\~no}}
\end{equation}

\textit{Donde:}
\begin{itemize}[noitemsep,leftmargin=2em]
    \item[$V_{total}$] = Volumen total del reactor (m³)
    \item[$Q_{medio}$] = Caudal medio diario de diseño (498.5 m³/d)
    \item[$DBO_{entrada}$] = DBO$_5$ del afluente (231 mg/L)
    \item[$COV_{dise\~no}$] = Carga orgánica volumétrica de diseño (1.00 kg DBO/m³·d)
\end{itemize}

Sustituyendo los valores del proyecto:

\begin{equation}
V_{total} = \frac{498.5 \times 231 / 1\,000}{1.00} = 115.1 \text{ m}^3
\end{equation}

El TRH resultante asociado a este volumen es:

\begin{equation}
TRH = \frac{V_{total}}{Q_{medio}} \times 24 = \frac{115.1}{498.5} \times 24 = 5.5 \text{ h}
\end{equation}

\textbf{[Ecuación 2] Volumen por compartimento:}

\begin{equation}
V_{comp} = \frac{V_{total}}{n} = \frac{115.1}{4} = 28.78 \text{ m}^3
\end{equation}

\textbf{[Ecuación 3] Dimensionamiento geométrico del módulo:}

El ancho del reactor se obtiene a partir del volumen de diseño, la profundidad útil y una relación geométrica objetivo $L_{comp}/W = 1.5$. Partiendo de la expresión del volumen total de un módulo:

\begin{equation}
V_{total} = n \times L_{comp} \times W \times H = n \times \left( \frac{L_{comp}}{W} \right) \times W^2 \times H
\end{equation}

Despejando el ancho $W$:

\begin{equation}
W = \sqrt{\frac{V_{total}}{n \times (L_{comp}/W)_{obj} \times H}} = \sqrt{\frac{115.1}{4 \times 1.5 \times 2.0}} = 1.79 \text{ m}
\end{equation}

El ancho teórico con un único módulo resultaría $W = 3.10$~m, superior al límite constructivo máximo de $6.0$~m. Por ello se divide la línea en $3$ módulos ABR paralelos, cada uno de $1.79$~m de ancho, con una separación libre entre ellos de $0.5$~m. El ancho total de la línea resulta $6.36$~m.

\textbf{[Ecuación 4] Dimensiones del reactor:}

A partir del ancho calculado, se obtienen las demás dimensiones geométricas del módulo:

\begin{equation}
L_{comp} = \frac{V_{comp}}{W \times H} = \frac{28.78}{1.79 \times 2.0} = 2.68 \text{ m} \quad \text{(largo por compartimento)}
\end{equation}

\begin{equation}
L_{total} = n \times L_{comp} = 4 \times 2.68 = 10.73 \text{ m} \quad \text{(largo total del módulo)}
\end{equation}

La velocidad ascensional real del módulo se verifica \textit{a posteriori} con el caudal que le corresponde:

\begin{equation}
v_{up,calc} = \frac{Q_{módulo}}{W \times H} = \frac{20.77 / 3}{1.79 \times 2.0} = 1.94 \text{ m/h}
\end{equation}


La profundidad total de construcción incluye la profundidad útil del líquido, la zona de acumulación de lodos y el bordo libre:

\begin{equation}
H_{total} = H + H_{lodos} + H_{bordo} = 2.0 + 0.30 + 0.30 = 2.60 \text{ m}
\end{equation}

\textbf{[Ecuación 5] Altura de abertura en los bafles (slot inferior):}

Cada pantalla o bafle interno no ocupa toda la altura del compartimento: deja una abertura en su parte inferior (slot) por la que el flujo debe pasar antes de ascender en el compartimento siguiente. La altura de este slot es un parámetro de diseño que determina la velocidad puntual del flujo en la transición entre compartimentos. Según OPS/OMS (2005) y Foxon et al. (2004), esta abertura se define como una fracción de la profundidad útil, típicamente entre el 15\,\% y el 25\,\% de $H$. Adoptando el valor central de la bibliografía ($f_{slot} = 0.20$):

\begin{equation}
h_{slot} = f_{slot} \times H = 0.20 \times 2.0 = 0.40 \text{ m}
\end{equation}

\textbf{[Ecuación 6] Velocidad en la abertura de los bafles:}

La velocidad del flujo al atravesar la abertura inferior de cada bafle (\textit{slot velocity}) es significativamente mayor que la velocidad ascensional media del compartimento, porque el área de paso queda reducida a $W \times h_{slot}$. Este parámetro es crítico: un valor excesivo genera turbulencia local en el fondo del compartimento siguiente, resuspende los lodos sedimentados y puede provocar su arrastre hacia el efluente. La bibliografía recomienda mantener esta velocidad por debajo de 1.0~m/h a caudal medio:

\begin{equation}
v_{slot} = \frac{Q_{módulo}}{W \times h_{slot}} = \frac{20.770 / 3}{1.79 \times 0.40} = 9.67 \text{ m/h}
\end{equation}

La carga orgánica volumétrica real del reactor resulta: $COV = 1.00$~kg\,DBO/(m³·d). 

\textbf{Limitaciones del método actual:}

El presente dimensionamiento cubre el cálculo del volumen total, número y tamaño de compartimentos, geometría del reactor, verificaciones hidráulicas y parámetros de carga. No cubre, en esta versión del cálculo: el modelado de la eficiencia de remoción de DBO en función del TRH (que requeriría un modelo cinético calibrado), el cálculo de la producción de biogás, el diseño del sistema de ventilación del biogás, el balance detallado de lodos en estado estacionario, ni el dimensionamiento de tuberías de entrada, salida y extracción de lodos.

\subsubsection{Verificación}

La subsección de verificación evalúa de forma sistemática si el reactor dimensionado satisface los criterios hidráulicos y geométricos de la bibliografía de referencia. Para cada verificación se describe primero el fenómeno físico que la motiva, después se desarrolla el cálculo con los valores del proyecto y, finalmente, se emite el juicio de cumplimiento. Se distinguen verificaciones \textbf{obligatorias} ---cuyo incumplimiento invalida el diseño--- de \textbf{complementarias} ---que orientan sobre la calidad del diseño sin invalidarlo--- y de \textbf{parámetros informativos} ---que no tienen criterio de cumplimiento en este nivel de diseño---.

\textbf{Verificaciones obligatorias:}

\textbf{[V-1] Rango bibliográfico de la velocidad ascensional de diseño}

El primer parámetro que debe justificarse es la propia elección de la velocidad ascensional de diseño $v_{up,dise\tilde{n}o} = 1.0$~m/h. Este valor no es un resultado del cálculo sino una decisión del proyectista: determina el área transversal del reactor (Ecuación~3) y, a través de ella, todas las demás dimensiones geométricas. Su validez depende de que se encuentre dentro del rango operativo aceptable reportado en la bibliografía técnica.

Las razones de este límite son físicas: velocidades demasiado altas generan shear suficiente para arrastrar los flóculos de lodo anaerobio hacia el efluente, empobreciendo progresivamente la biomasa activa del reactor; velocidades demasiado bajas producen reactores sobredimensionados sin mejora proporcional en la eficiencia biológica, pues el limitante del proceso pasa a ser la cinética de la metanogénesis y no el tiempo de contacto. Según Foxon et al. (2004) y Barber \& Stuckey (1999), el rango operativo aceptable para aguas residuales domésticas a temperatura ambiente es:

\begin{equation}
0.2 \text{ m/h} \leq v_{up,dise\tilde{n}o} \leq 1.5 \text{ m/h}
\end{equation}

El valor adoptado en este proyecto es $v_{up,dise\tilde{n}o} = 1.0$~m/h.

\textit{Nota importante sobre la verificación de la velocidad a caudal medio:} el ancho $W$ ya no se calcula imponiendo la velocidad ascensional de diseño, sino a partir de criterios geométricos (relación $L_{comp}/W$ y profundidad $H$). Por tanto, la velocidad ascensional calculada a caudal medio ($v_{up,calc} = Q_{módulo} / (W \times H)$) es ahora un resultado real del diseño, no una identidad tautológica. Esto permite verificar objetivamente si el reactor opera dentro del rango admisible. La verificación más crítica sigue siendo la del caudal pico (V-2).

Estado V-1: \textbf{CUMPLE} ($0.2 \leq 1.0 \leq 1.5$~m/h)

\textbf{[V-2] Velocidad ascensional a caudal máximo horario}

Esta es la verificación hidráulica más crítica del diseño. Durante los picos de consumo ---generalmente en las mañanas y mediodías en sistemas domésticos--- el caudal supera al medio en un factor de pico $f_p$. Si en esas condiciones la velocidad ascensional rebasa el umbral de arrastre, la biomasa abandona el reactor con el efluente y el rendimiento de depuración se deteriora progresivamente. A diferencia del UASB, el ABR carece de separador trifásico que retenga los sólidos en esas condiciones: la única salvaguarda es que la geometría garantice una velocidad admisible incluso en pico.

El caudal pico horario se obtiene aplicando el factor de pico al caudal medio de diseño:

\begin{equation}
Q_{max} = f_p \times Q_{medio} = 2.8 \times 20.770 = 57.540 \text{ m}^3\text{/h}
\end{equation}

Con este caudal, la velocidad ascensional en pico resulta:

\begin{equation}
v_{up,max} = \frac{Q_{max}}{W \times H} = \frac{57.540}{1.79 \times 2.0} = 5.36 \text{ m/h}
\end{equation}

El límite bibliográfico para caudal pico, según Barber \& Stuckey (1999), es de 2.0~m/h. Por encima de este umbral se ha observado exportación significativa de sólidos hacia el efluente en reactores piloto y a escala real.

Criterio: $v_{up,max} \leq 2.0$~m/h

Estado V-2: \textbf{NO CUMPLE} (5.36~m/h $>$ 2.0~m/h)

\textbf{[V-3] Verificación del tiempo de retención hidráulico por geometría}

El TRH de diseño fija el volumen total del reactor (Ecuación~1). Sin embargo, las dimensiones finales resultan de una cadena de operaciones aritméticas ---división por área, multiplicación por número de compartimentos--- que en la práctica pueden ajustarse a valores constructivos enteros. Esta verificación confirma que el volumen efectivo de la geometría adoptada reproduce, o supera, el TRH de diseño. Si el TRH calculado resultara inferior al de diseño, significaría que las dimensiones son insuficientes y que el tiempo de contacto nominal no se alcanzaría en la práctica.

\begin{equation}
TRH_{calc} = \frac{L_{total} \times W \times H_{util}}{Q_{medio}} \times 24 = \frac{10.73 \times 1.79 \times 2.0}{498.5} \times 24 = 5.5 \text{ h}
\end{equation}

Criterio: $TRH_{calc} \geq 6$~h

Estado V-3: \textbf{NO CUMPLE} (5.5~h $<$ 6~h)

\textbf{[V-4] Largo mínimo de compartimento}

El largo de cada compartimento ($L_{comp}$) debe satisfacer dos criterios independientes que responden a razones diferentes. El criterio constructivo establece un mínimo absoluto de 1.0~m: por debajo de esta dimensión el compartimento es impracticable desde el punto de vista de la ejecución de obra, la inspección y la extracción de lodos. El criterio hidráulico-bibliográfico, según OPS/OMS (2005), establece que $L_{comp} \geq H_{util}$: esta condición garantiza que el recorrido horizontal del flujo sea suficiente para que los sólidos gruesos no degradados sedimenten antes de alcanzar el bafle siguiente, evitando cortocircuitos hidráulicos. Ambos criterios deben cumplirse simultáneamente, por lo que el largo mínimo efectivo es el mayor de los dos valores:

\begin{equation}
L_{comp,min} = \max\!\left(H_{util},\; 1.0\right) = \max\!\left(2.0,\; 1.0\right) = 2.0 \text{ m}
\end{equation}

Valor calculado: $L_{comp} = 2.68$~m

Estado V-4: \textbf{CUMPLE} ($L_{comp} = 2.68$~m $\geq$ 2.0~m)

\textbf{[V-ancho] Ancho del módulo ABR}

El ancho del módulo es una verificación constructiva que garantiza que cada unidad ABR puede edificarse y operarse sin alcanzar dimensiones excesivas. Un ancho superior al máximo admisible dificulta la ejecución de obra, la colocación de pantallas internas y la distribución uniforme del afluente en el plano transversal. Si el ancho teórico calculado para un único módulo excede el límite, la línea se divide en módulos paralelos hasta satisfacer el criterio.

Criterio: $W \leq 6.0$~m.

Valor calculado: $W = 1.79$~m (con $3$ módulo(s) paralelo(s)).

Estado V-ancho: \textbf{CUMPLE} (1.79~m $\leq$ 6.0~m)

\textbf{[V-largo] Largo total del módulo ABR}

El largo total del módulo es una verificación constructiva complementaria al control del ancho. Un módulo excesivamente largo dificulta la ejecución de obra, la distribución uniforme del afluente a lo largo del reactor y el mantenimiento de las pantallas internas. Si el largo calculado para un único módulo supera el límite admisible, se incrementa el número de módulos paralelos hasta que ambas dimensiones (ancho y largo) sean constructivamente razonables.

Criterio: $L_{total} \leq 12.0$~m.

Valor calculado: $L_{total} = 10.73$~m (con $3$ módulo(s) paralelo(s)).

Estado V-largo: \textbf{CUMPLE} (10.73~m $\leq$ 12.0~m)

\textbf{[V-5] Velocidad en la abertura inferior de los bafles (\textit{slot velocity})}

Cada pantalla interna del ABR deja una abertura en su parte inferior (denominada \textit{slot} en la bibliografía anglosajona) a través de la cual el flujo pasa de un compartimento al siguiente antes de ascender. Dado que el área de esta abertura ($W \times h_{slot}$) es sustancialmente menor que el área transversal total del compartimento ($W \times H$), la velocidad puntual del flujo en el slot es notablemente mayor que la velocidad ascensional media.

Esta aceleración local tiene dos efectos negativos cuando supera el umbral admisible: primero, genera turbulencia en el fondo del compartimento receptor, resuspendiendo los flóculos de lodo ya sedimentados; segundo, puede crear zonas de mezcla intensa que cortocircuitan la separación funcional entre compartimentos, que es uno de los fundamentos del buen rendimiento del ABR. OPS/OMS (2005) y Foxon et al. (2004) recomiendan definir la abertura como el 20\,\% de la profundidad útil ($f_{slot} = 0.20$) y mantener la velocidad en el slot por debajo de 1.0~m/h a caudal medio:

\begin{equation}
h_{slot} = f_{slot} \times H_{util} = 0.20 \times 2.0 = 0.40 \text{ m}
\end{equation}

La velocidad en la abertura a caudal medio resulta:

\begin{equation}
v_{slot} = \frac{Q_{módulo}}{W \times h_{slot}} = \frac{20.770 / 3}{1.79 \times 0.40} = 9.67 \text{ m/h}
\end{equation}

A modo informativo, a caudal máximo por módulo:

\begin{equation}
v_{slot,max} = \frac{Q_{max,módulo}}{W \times h_{slot}} = \frac{57.540 / 3}{1.79 \times 0.40} = 26.79 \text{ m/h}
\end{equation}

Criterio: $v_{slot} \leq 1.0$~m/h (caudal medio)

Estado V-5: \textbf{NO CUMPLE} (9.67~m/h $>$ 1.0~m/h)

\textbf{[V-6] Relación profundidad útil / ancho del reactor ($H/W$)}

La relación $H/W$ describe la forma de la sección transversal de cada compartimento. Su importancia es reconocida en la bibliografía reciente (Sinha et al., 2009): valores extremos en cualquier sentido comprometen la uniformidad del flujo vertical ascendente. Un reactor excesivamente ancho en relación con su profundidad ($H/W < 0.5$) tiende a tener distribución no uniforme del flujo en el plano transversal, con zonas muertas en los extremos laterales donde la velocidad ascensional es insuficiente para mantener el manto de lodos en suspensión activa. Por el contrario, un reactor excesivamente estrecho ($H/W > 2.0$) dificulta la distribución uniforme del afluente en el ancho y puede crear flujos preferenciales en la zona central que reducen el contacto efectivo entre el sustrato y la biomasa. El rango recomendado es:

\begin{equation}
0.5 \leq \frac{H_{util}}{W} \leq 2.0
\end{equation}

El valor calculado para este diseño es:

\begin{equation}
\frac{H_{util}}{W} = \frac{2.0}{1.79} = 1.12
\end{equation}

Estado V-6: \textbf{CUMPLE} ($0.5 \leq 1.12 \leq 2.0$)

\textbf{[V-COV] Carga orgánica volumétrica (COV)}

La carga orgánica volumétrica expresa la cantidad de materia orgánica que ingresa por unidad de volumen y por día. Es el parámetro de referencia más utilizado en la bibliografía de reactores anaerobios para comparar diseños y evaluar si el reactor opera en un régimen compatible con el mantenimiento de una biomasa estable. Un valor demasiado bajo indica un reactor sobredimensionado desde el punto de vista biológico, lo que no mejora la calidad del efluente proporcionalmente al costo adicional y puede generar problemas de arranque o inestabilidad del manto de lodos.

La DBO se expresa en mg/L, unidad equivalente a g/m³; dividirla entre 1\,000 convierte g a kg. La expresión completa con sus unidades explícitas es:

\begin{equation}
COV = \frac{Q_{medio} \,[\text{m}^3/\text{d}] \times DBO_{entrada} \,[\text{g/m}^3] \,/\, 1\,000}{V_{total} \,[\text{m}^3]} = \frac{498.5 \times 231 / 1\,000}{115.1} = 1.00 \text{ kg\,DBO/(m}^3\!\cdot\!\text{d)}
\end{equation}

Para este proyecto, la carga orgánica máxima admisible se corrige por temperatura mediante un factor de Arrhenius, dado que el reactor opera en condiciones de clima cálido:

\begin{equation}
f_T = \theta^{T - 20} = 1.04^{25.6 - 20} = 1.246
\end{equation}

\begin{equation}
COV_{m\acute{a}x} = 1.0 \times f_T = 1.0 \times 1.246 = 1.25 \text{ kg\,DBO/(m}^3\!\cdot\!\text{d)}
\end{equation}

Criterio obligatorio: $0.50 \leq COV \leq 1.25$~kg\,DBO/(m³·d).

Estado V-COV: \textbf{CUMPLE} (1.00~kg\,DBO/(m³·d) dentro del rango ajustado por temperatura)

\textbf{[V-TRHmin] TRH mínimo del proceso anaerobio}

Cuando el reactor no se dimensiona directamente por TRH, el tiempo de retención hidráulico resultante debe verificarse contra un mínimo absoluto para garantizar que el proceso anaerobio dispone de tiempo suficiente de contacto y que la biomasa no es arrastrada por washout. Valores de TRH inferiores a 6.0~h son considerados inadecuados para reactores anaerobios domésticos sin aporte de calor o control de temperatura.

Criterio obligatorio: $TRH \geq 6.0$~h.

Estado V-TRHmin: \textbf{NO CUMPLE - riesgo de washout} (5.5~h $<$ 6.0~h)

\textbf{Verificaciones complementarias:}

\textbf{[V-7] Relación largo / ancho del compartimento ($L_{comp} / W$)}

La relación $L_{comp} / W$ describe la esbeltez en planta de cada compartimento. Compartimentos muy alargados ($L_{comp} / W > 2.0$) pueden presentar distribución no uniforme del afluente a lo largo del ancho; compartimentos muy cuadrados o anchos ($L_{comp} / W < 0.5$) tienden hacia un comportamiento de mezcla completa en el plano horizontal, reduciendo la eficiencia de la separación en serie entre compartimentos. El rango orientativo recomendado por la OPS/OMS (2005) es:

\begin{equation}
0.5 \leq \frac{L_{comp}}{W} \leq 2.0
\end{equation}

Valor calculado: $L_{comp} / W = 2.68 / 1.79 = 1.50$

Estado V-7: \textbf{DENTRO RANGO} (verificación complementaria)

\textbf{[V-8] Número mínimo de compartimentos}

El número de compartimentos determina el grado de separación funcional entre las fases del proceso anaerobio. Con un solo compartimento el ABR no tiene ventaja sobre un digestor de mezcla completa. A partir de cuatro compartimentos se logra una separación razonablemente clara entre las comunidades hidrolítico-acidogénicas (primer y segundo compartimento) y las metanogénicas (últimos compartimentos). La OPS/OMS (2005) establece un mínimo de 3 compartimentos para aplicaciones de saneamiento doméstico descentralizado:

\begin{equation}
n \geq 3
\end{equation}

Valor adoptado: $n = 4$

Estado V-8: \textbf{CUMPLE}

\textbf{Parámetros informativos:}

\textbf{[V-9] Carga orgánica volumétrica (COV)}

La carga orgánica volumétrica expresa la cantidad de materia orgánica que ingresa por unidad de volumen y por día. Es el parámetro de referencia más utilizado en la bibliografía de reactores anaerobios para comparar diseños y evaluar si el reactor opera en un régimen compatible con el mantenimiento de una biomasa estable. No constituye un criterio obligatorio en este diseño, pero permite situar el proyecto en el contexto de experiencias reportadas.

La DBO se expresa en mg/L, unidad equivalente a g/m³; dividirla entre 1\,000 convierte g a kg. La expresión completa con sus unidades explícitas es:

\begin{equation}
COV = \frac{Q_{medio} \,[\text{m}^3/\text{d}] \times DBO_{entrada} \,[\text{g/m}^3] \,/\, 1\,000}{V_{total} \,[\text{m}^3]} = \frac{498.5 \times 231 / 1\,000}{115.1} = 1.00 \text{ kg\,DBO/(m}^3\!\cdot\!\text{d)}
\end{equation}

Rango bibliográfico para ABR doméstico a temperatura ambiente: 0.5--5.0~kg\,DBO/(m³·d) (Barber \& Stuckey, 1999; OPS/OMS, 2005). Nota: en este diseño la COV también constituye una verificación obligatoria (V-COV).

\textbf{[V-10] Estimación de la producción de lodos y frecuencia de purga}

La zona de acumulación de lodos reservada en el dimensionamiento ($H_{lodos} = 0.30$~m) debe ser coherente con la tasa real de producción de biomasa anaerobia en el reactor. Sin esta comprobación, el espacio reservado para lodos es solo un criterio empírico desconectado de las condiciones reales del proyecto. La OPS/OMS (2005, Capítulo~6) proporciona la siguiente expresión para una estimación de orden de magnitud:

\begin{equation}
P_{lodos} = Y_{obs} \times DBO_{rem}
\end{equation}

Donde $Y_{obs} \approx 0.10$~kg\,SSV/kg\,DBO\,removida es el rendimiento observado en digestión anaerobia mesofílica de aguas residuales domésticas (Rittmann \& McCarty, 2001), y $DBO_{rem}$ es la DBO removida diariamente. Adoptando la eficiencia de referencia bibliográfica del 75\,\%:

\begin{equation}
DBO_{rem} = \eta_{ref} \times DBO_{entrada} \times \frac{Q_{medio}}{1\,000} = 0.75 \times 231 \times \frac{498.5}{1\,000} = 86.35 \text{ kg\,DBO/d}
\end{equation}

\begin{equation}
P_{lodos} = 0.10 \times 86.35 = 8.63 \text{ kg\,SSV/d}
\end{equation}

Adoptando una densidad de lodos húmedos anaerobios espesados de $\rho = 1\,040$~kg/m³, el volumen acumulado anualmente es:

\begin{equation}
V_{lodos,año} = \frac{P_{lodos} \times 365}{\rho} = \frac{8.63 \times 365}{1\,040} = 3.03 \text{ m}^3\text{/año}
\end{equation}

El volumen disponible en la zona de lodos del reactor es:

\begin{equation}
V_{zona\,lodos} = L_{total} \times W \times H_{lodos} = 10.73 \times 1.79 \times 0.30 = 5.76 \text{ m}^3
\end{equation}

Esto implica una frecuencia indicativa de purga de lodos de aproximadamente \textbf{18 meses}, asumiendo extracción cuando la zona de acumulación alcanza el 80\,\% de su capacidad. Este valor es una estimación de orden de magnitud y deberá ajustarse en la puesta en marcha del sistema en función de las mediciones reales de producción de lodos.

\textbf{Resumen de verificaciones:}

\begin{table}[H]
\centering
\caption{Resumen de verificaciones del reactor ABR / RAP}
\begin{tabular}{llp{5.5cm}}
\toprule
\textbf{Verificación} & \textbf{Estado} & \textbf{Valor calculado vs. criterio} \\
\midrule
V-1 Rango v\_up diseño & \checkmark & 1.0 m/h $\in$ [0.2--1.5] m/h \\
V-2 v\_up caudal máximo & \texttimes & 5.36 m/h $\leq$ 2.0 m/h \\
V-3 TRH verificado & \texttimes & 5.5 h $\geq$ 6 h \\
V-4 Largo mínimo compartimento & \checkmark & L\_comp = 2.68 m $\geq$ 2.0 m \\
V-5 Velocidad abertura bafle & \texttimes & 9.67 m/h $\leq$ 1.0 m/h \\
V-6 Relación H/W & \checkmark & 1.12 $\in$ [0.5--2.0] \\
V-COV Carga orgánica volumétrica & \checkmark & 1.00 kg DBO/m³·d $\geq$ 0.50 kg DBO/m³·d \\
V-ancho Ancho del módulo & \checkmark & 1.79 m $\leq$ 6.00 m \\
V-largo Largo total del módulo & \checkmark & 10.73 m $\leq$ 12.00 m \\
V-TRHmin TRH mínimo proceso & \texttimes & 5.5 h $\geq$ 6.0 h \\
\bottomrule
\end{tabular}
\end{table}

\textbf{ATENCIÓN:} las siguientes verificaciones NO cumplen y requieren revisión del diseño: V-2 v\_up caudal máximo, V-3 TRH verificado, V-5 Velocidad abertura bafle, V-TRHmin TRH mínimo proceso.

\textit{Notas: (a) El proyecto NO calcula la eficiencia de remoción de DBO a partir de parámetros operacionales. Las eficiencias del 70--95\,\% reportadas en la bibliografía para ABR con TRH de 6~h son referencias y no predicciones garantizadas. (b) La estimación de producción de lodos (V-10) adopta una eficiencia de referencia del 75\,\%; el valor real dependerá de las condiciones de operación.}

\subsubsection{Resultados}

\begin{figure}[H]
\centering
\includegraphics[width=0.95\textwidth]{figuras/Esquema_ABR.png}
\caption{Esquema del Reactor Anaerobio con Pantallas (ABR / RAP) -- vista longitudinal}
\label{fig:esquema_abr}
\end{figure}

El dimensionamiento del Reactor Anaerobio con Pantallas (ABR / RAP) ha producido las siguientes dimensiones y parámetros de operación:

El esquema muestra la configuración típica del ABR con flujo en zigzag (\textit{up-flow} y \textit{down-flow}) a través de los compartimentos, las pantallas internas (\textit{baffles}) que dirigen el flujo, las zonas de lodos, líquida y biogás, así como los puntos de entrada, salida y ventilación.

\begin{longtable}{p{6cm}p{3cm}p{5cm}}
\caption{Resultados del dimensionamiento ABR / RAP} \\
\toprule
\textbf{Parámetro} & \textbf{Valor} & \textbf{Observación} \\
\midrule
\endfirsthead
\multicolumn{3}{c}{
    {\bfseries \tablename\ \thetable{} -- continuación de la página anterior}
} \\
\toprule
\textbf{Parámetro} & \textbf{Valor} & \textbf{Observación} \\
\midrule
\endhead
\midrule
\multicolumn{3}{r}{Continúa en la siguiente página...} \\
\endfoot
\bottomrule
\endlastfoot

\multicolumn{3}{l}{\textit{Parámetros de entrada:}} \\
Caudal medio ($Q_{medio}$) & 498.5 m³/d & Caudal por línea de tratamiento \\
Caudal máximo ($Q_{max}$) & 1380.9 m³/d & Factor pico: 2.8 \\
DBO afluente & 231 mg/L & Caracterización afluente \\
\midrule

\multicolumn{3}{l}{\textit{Criterios de diseño adoptados:}} \\
Tiempo retención hidráulico (TRH) & 6 h & Criterio del proyecto \\
Número de compartimentos ($n$) & 4 & 4--8 recomendado \\
Velocidad ascensional diseño & 1.0 m/h & $\leq$ 2.0 m/h límite pico \\
Profundidad útil ($H$) & 2.0 m & 1.5--3.0 m rango \\
\midrule

\multicolumn{3}{l}{\textit{Dimensiones geométricas:}} \\
Volumen total ($V_{total}$) & 115.1 m³ & Calculado por COV \\
Volumen por compartimento & 28.78 m³ & Distribución uniforme \\
Número de módulos paralelos & 3 & Por línea de tratamiento \\
Área sección transversal & 3.58 m² & Geometría real del módulo \\
Ancho del módulo ($W$) & 1.79 m & $\leq$ 6.0 m \\
Ancho total de la línea & 6.36 m & Incluye separación entre módulos \\
Largo por compartimento & 2.68 m & $L_{comp} \geq H$ verificado \\
Largo total del módulo ($L_{total}$) & 10.73 m & $n \times L_{comp}$ \\
Profundidad total construcción & 2.60 m & $H + H_{lodos} + H_{bordo}$ \\
Área módulo individual & 19.19 m² & $L_{total} \times W_{módulo}$ \\
Área total línea en planta & 68.30 m² & Huella constructiva total \\
\midrule

\multicolumn{3}{l}{\textit{Parámetros hidráulicos:}} \\
Velocidad ascensional (caudal medio) & 1.94 m/h & NO CUMPLE \\
Velocidad ascensional (caudal máximo) & 5.36 m/h & NO CUMPLE \\
TRH calculado & 5.5 h & NO CUMPLE \\
Carga hidráulica superficial (CHS) & 7.30 m/d & Parámetro informativo \\
\midrule

\multicolumn{3}{l}{\textit{Parámetros de proceso (informativos):}} \\
Carga orgánica volumétrica (COV) & 1.00 kg/m³·d & Rango ref: 0.5--5.0 \\
\midrule

\multicolumn{3}{l}{\textit{Estado de verificaciones:}} \\
Verificación velocidad (medio) & NO CUMPLE & Criterio: $\leq$ 1.0 m/h \\
Verificación velocidad (máx) & NO CUMPLE & Criterio: $\leq$ 2.0 m/h \\
Verificación TRH & NO CUMPLE & Criterio: $\geq$ 6 h \\
Verificación largo comp. & CUMPLE & Criterio: $L_{comp} \geq H$ \\
Verificación ancho módulo & CUMPLE & Criterio: $W \leq$ 6.0 m \\
Verificación largo total & True & Criterio: $L_{total} \leq$ 12.0 m \\
Verificación COV & True & Criterio: $\geq$ 0.50 kg/m³·d \\
Verificación TRH mínimo & False & Criterio: $\geq$ 6.0 h \\
\bottomrule
\multicolumn{3}{l}{\textbf{Estado general del diseño:} \textbf{REQUIERE REVISIÓN}} \\
\end{longtable}

\textbf{Notas sobre los resultados:}

El reactor ABR / RAP dimensionado consta de 4 compartimentos idénticos con una geometría rectangular de 10.73~m de largo total, 1.79~m de ancho y 2.60~m de profundidad total de construcción. La profundidad útil del líquido es de 2.0~m, con una zona adicional de 0.30~m reservada para acumulación de lodos y 0.30~m de bordo libre.

El diseño opera con una velocidad ascensional de 1.94~m/h a caudal medio, valor que se eleva hasta 5.36~m/h durante los picos de caudal. La velocidad a caudal medio cumple con el límite operativo, pero la velocidad de pico supera el umbral bibliográfico admisible (2.0~m/h), lo que indica que el reactor puede experimentar arrastre de biomasa en condiciones de caudal máximo y debe ser revisado.

El tiempo de retención hidráulico resultante es de 5.5~h (0.23~días). Este valor proporciona el tiempo de contacto necesario entre el agua residual y la biomasa anaerobia para la degradación de la materia orgánica.

\textbf{Alcance y limitaciones declaradas:}

Este dimensionamiento cubre el volumen del reactor, la geometría básica y las verificaciones hidráulicas fundamentales. No incluye: (a) predicción de la eficiencia de remoción de DBO o DQO mediante modelos cinéticos; (b) cálculo de la producción de biogás ni dimensionamiento del sistema de ventilación; (c) balance detallado de producción de lodos y frecuencia de extracción; (d) diseño de las pantallas internas (altura, espesor, configuración de aberturas); (e) dimensionamiento de tuberías de entrada, salida y extracción de lodos. Estos elementos constituyen extensiones futuras del diseño que deben desarrollarse según las condiciones específicas del proyecto.

\subsection{Biofiltro de Carga Mecanizada e Hidráulica}

\subsubsection{Dimensionamiento}

El biofiltro percolador de carga mecanizada hidráulica, conocido internacionalmente como Trickling Filter (TF) o biofiltro de goteo, constituye uno de los procesos biológicos más antiguos y probados para el tratamiento secundario de aguas residuales. Su principio de funcionamiento se basa en el contacto intermitente entre el agua residual y una biopelícula aerobia adherida a un medio de relleno estructurado, permitiendo la oxidación de la materia orgánica biodegradable mediante procesos biológicos aerobios.

El término ``carga mecanizada hidráulica'' hace referencia a que el afluente se aplica sobre el medio mediante un sistema de distribución rotatorio accionado hidráulicamente, en contraste con los sistemas de alimentación por gravedad simple. Esta configuración garantiza una distribución uniforme del agua residual sobre toda la superficie del medio, promoviendo el desarrollo homogéneo de la biopelícula y maximizando la eficiencia de transferencia de oxígeno y substrato.

\subsubsection*{Fundamentos del Proceso y Desarrollo de la Biopelícula}

El biofiltro opera como un reactor biológico de película fija (\textit{attached growth}) donde los microorganismos forman una capa biológica (biopelícula o \textit{biofilm}) de aproximadamente 0.1 a 0.5 mm de espesor sobre la superficie del medio de relleno. Esta biopelícula está compuesta por una comunidad microbiana compleja que incluye bacterias aerobias, hongos, protozoos y metazoos (principalmente rotíferos y nemátodos), organizada en capas funcionales distintas.

La capa externa de la biopelícula, en contacto directo con el aire que circula por el medio, mantiene condiciones estrictamente aerobias y es donde ocurre la oxidación de la materia orgánica. La capa interna, adyacente a la superficie del medio, puede desarrollar condiciones anaerobias o anóxicas, particularmente cuando el espesor de la biopelícula excede los 0.3 mm. El aire circula por el biofiltro mediante convección natural, impulsado por la diferencia de temperatura entre el aire ambiente y el interior del filtro, proporcionando el oxígeno necesario para mantener la actividad aerobia de los microorganismos.

El agua residual distribuida sobre el medio percola gravitacionalmente a través de los intersticios del relleno, entrando en contacto con la biopelícula. Durante este contacto, los contaminantes orgánicos se adsorben en la superficie de la biopelícula y son metabolizados por los microorganismos. Simultáneamente, el oxígeno disuelto en el afluente y el transferido desde la fase gaseosa es consumido en la respiración aerobia de la biomasa. El exceso de biomasa producida (lodo biológico o \textit{sloughings}) se desprende periódicamente de la superficie del medio por efecto del arrastre hidráulico, siendo recolectado en el sistema de drenaje inferior y requiriendo separación en un sedimentador secundario.

\subsubsection*{Selección de la Metodología de Diseño}

El manual técnico establece dos rutas de diseño fundamentales para el dimensionamiento de biofiltros percoladores, cada una aplicable a condiciones específicas de carga orgánica, disponibilidad de energía y requisitos de efluente:

\textbf{Ruta A -- Carga Hidráulica Convencional (modelo NRC):} Esta metodología se aplica cuando la carga orgánica superficial es moderada, típicamente para cargas orgánicas volumétricas menores o iguales a 0.40 kg DBO/m³·d. La Ruta A no utiliza recirculación del efluente tratado, operando en flujo simple (\textit{once-through}). El modelo de diseño es el desarrollado por el National Research Council (NRC) de Estados Unidos, derivado del análisis estadístico de datos operacionales de instalaciones militares durante la Segunda Guerra Mundial. Este modelo es conservador y ha demostrado su eficacia en instalaciones con cargas moderadas y requisitos de efluente estándar.

\textbf{Ruta B -- Carga Mecanizada de Alta Carga (modelo Germain):} Esta metodología se selecciona cuando la carga orgánica superficial excede el umbral de 0.40 kg DBO/m³·d, cuando se requiere una eficiencia de remoción superior, o cuando se dispone de energía para bombeo de recirculación. La Ruta B incorpora la recirculación del efluente tratado hacia la entrada del biofiltro, con relaciones de recirculación típicas entre $R = 0.5$ y $R = 4.0$. El modelo de diseño es el propuesto por Germain (1966) para medios plásticos estructurados, que describe la cinética de remoción mediante una ecuación de primer orden modificada que incorpora el efecto de la recirculación sobre la concentración afluente.

\begin{table}[H]
\centering
\caption{Criterios de selección entre Rutas de diseño para biofiltros TAF}
\label{tab:seleccion_ruta_taf}
\begin{tabular}{p{4.5cm}p{5cm}p{5cm}}
\toprule
\textbf{Criterio} & \textbf{Ruta A (NRC)} & \textbf{Ruta B (Germain)} \\
\midrule
Carga orgánica volumétrica (COS) & $\leq 0.40$ kg/m³·d & $> 0.40$ kg/m³·d \\
Recirculación & No (flujo simple) & Sí ($R = 0.5$--$4.0$) \\
Requerimiento energético & Mínimo (solo distribución) & Moderado (bombeo recirculación) \\
Eficiencia típica DBO$_5$ & 60--80\% & 70--90\% (con $R$ óptimo) \\
Modelo matemático & NRC (empírico) & Germain (cinético) \\
Aplicación típica & Cargas moderadas, área disponible & Alta carga, estricto efluente \\
\bottomrule
\end{tabular}
\end{table}

El algoritmo de selección evalúa la carga orgánica preliminar estimada para determinar la ruta más adecuada al proyecto:

\begin{equation}
\text{COS}_{\text{prelim}} = \frac{W}{V_{\text{estimado}}} = \frac{Q \cdot S_0}{V_{\text{estimado}}}
\end{equation}

Para este proyecto, Se selecciona la Ruta B (Carga Mecanizada de Alta Carga) según el criterio: COS_prelim = 0.60 kg/m³·d > 0.40 → Ruta B. Esta ruta utiliza el modelo Germain con recirculación R=1.0, apropiada para alta carga orgánica y mejora la eficiencia de tratamiento. El criterio de selección adoptado garantiza que la metodología de diseño es coherente con las condiciones de carga y los objetivos de calidad del efluente establecidos.

\subsubsection*{Parámetros de Diseño y Rangos de Operación}

El diseño del biofiltro TAF se fundamenta en parámetros que han sido establecidos a través de décadas de experiencia operacional en instalaciones a escala real. La Tabla~\ref{tab:parametros_taf} consolida los parámetros principales con sus rangos típicos y los valores adoptados para este proyecto.

\begin{table}[H]
\centering
\caption{Parámetros de diseño para biofiltro TAF -- Ruta B}
\label{tab:parametros_taf}
\begin{tabular}{p{4cm}ccc}
\toprule
\textbf{Parámetro} & \textbf{Rango típico} & \textbf{Valor adoptado} & \textbf{Fuente} \\
\midrule
Profundidad medio $D$ & 2.0--8.0 m & 3.00 m & Metcalf \& Eddy \\
Constante cinética $k_{20}$ & 0.06--0.10 m/h & 0.068 m/h & Germain (1966) \\
Coeficiente temperatura $\theta$ & 1.03--1.06 & 1.035 & Arrhenius modificado \\
Exponente modelo $n$ (plástico) & 0.50 & 0.5 & WEF MOP-8 \\
Carga hidráulica superficial CHS & 0.5--4.0 m³/m²·h & 4.330 m³/m²·h & Proyecto \\
\bottomrule
\end{tabular}
\end{table}

La profundidad del medio de 3.00 m seleccionada se encuentra dentro del rango recomendado para biofiltros con medio plástico estructurado. Profundidades mayores aumentan el tiempo de contacto y la eficiencia de remoción, pero incrementan la carga estructural sobre el medio y el costo de construcción. La profundidad adoptada representa un compromiso óptimo entre eficiencia de tratamiento y viabilidad constructiva.

La constante cinética $k_{20} = 0.068$ m/h corresponde al valor recomendado por Metcalf \& Eddy (2014) para medios plásticos aleatorios modernos con superficie específica de aproximadamente 100 m²/m³. Este valor ha sido validado en múltiples instalaciones y representa una estimación conservadora para condiciones de operación normales.

\subsubsection*{Caracterización del Afluente}

El biofiltro recibe el efluente del tratamiento primario (UASB o sedimentador primario), con una carga orgánica residual que requiere remoción biológica aerobia. La Tabla~\ref{tab:entrada_taf} presenta los parámetros de entrada característicos del afluente al biofiltro.

\begin{table}[H]
\centering
\caption{Caracterización del afluente -- Biofiltro TAF}
\label{tab:entrada_taf}
\begin{tabular}{llll}
\toprule
\textbf{Parámetro} & \textbf{Símbolo} & \textbf{Valor} & \textbf{Unidad} \\
\midrule
Caudal promedio diario & $Q$ & 498.5 & m³/d \\
Caudal máximo horario & $Q_{\text{max}}$ & 1380.9 & m³/d \\
Factor de pico & -- & 2.8 & -- \\
DBO₅ afluente & $S_0$ & 57.7 & mg/L \\
SST afluente & $\text{SST}_0$ & 62.4 & mg/L \\
Temperatura del agua & $T$ & 25.6 & °C \\

\bottomrule
\end{tabular}
\end{table}

\subsubsection*{Carga Orgánica Afluente}

La carga orgánica afluente al biofiltro, calculada como el producto del caudal por la concentración de DBO₅, resulta:

\begin{equation}
W = Q \cdot S_0 = 498.5\ \text{m}^3\text{/d} \times 57.7\ \text{g/m}^3 \times 10^{-3} = 28.78\ \text{kg DBO/d}
\end{equation}
\captionequation{Carga orgánica diaria afluente al biofiltro TAF}

Esta carga orgánica representa la masa de contaminante biodegradable que debe ser removida por la biopelícula aerobia desarrollada en el medio filtrante.

\subsubsection*{Modelo de Diseño Germain -- Ruta B (con recirculación)}

Para la Ruta B se utiliza el modelo cinético desarrollado por Germain (1966) para biofiltros con medio plástico estructurado y recirculación. Este modelo describe la remoción de DBO$_5$ mediante una cinética de primer orden que incorpora el efecto de la recirculación sobre la concentración afluente diluida y la carga hidráulica total aplicada.

El modelo Germain parte de la observación experimental de que la remoción de substrato en biofiltros con medio plástico sigue una cinética proporcional a la concentración remanente, pero modulada por la tasa hidráulica. La forma matemática del modelo es:

\begin{equation}
\frac{S_e}{S_0'} = \exp\!\left(-\frac{k_T \cdot D}{Q_T^n}\right)
\end{equation}
\captionequation{Modelo de Germain para biofiltro con recirculación -- Ruta B}

Donde $S_e$ es la concentración de DBO$_5$ en el efluente, $S_0'$ es la concentración afluente diluida por la recirculación, $k_T$ es la constante cinética corregida por temperatura, $D$ es la profundidad del medio, $Q_T$ es la carga hidráulica superficial total (con recirculación), y $n$ es un exponente empírico que depende del tipo de medio ($n = 0.50$ para medio plástico).

La recirculación cumple tres funciones fundamentales en el biofiltro: (1) diluye la concentración de DBO$_5$ del afluente crudo, reduciendo la carga orgánica puntual y previniendo la sobrecarga del sistema; (2) aumenta la carga hidráulica superficial, promoviendo un riego más uniforme de la biopelícula y previniendo su desecación durante períodos de bajo caudal; y (3) mejora la distribución del flujo a través del medio, reduciendo la formación de canales preferenciales.

La concentración afluente diluida por efecto de la recirculación se calcula como:

\begin{equation}
S_0' = \frac{S_0}{1 + R} = \frac{57.7}{1 + 1.0} = 28.9\ \text{mg/L}
\end{equation}
\captionequation{Concentración afluente diluida por recirculación}

La relación de recirculación adoptada $R = 1.0$ significa que por cada unidad de caudal afluente, se recircula 1.0 unidades de efluente tratado, resultando en un caudal total aplicado al filtro de:

\begin{equation}
Q_T = Q \cdot (1 + R) = 20.77 \times 2.0 = 41.54\ \text{m}^3\text{/h} = 498.5\ \text{m}^3\text{/d}
\end{equation}
\captionequation{Caudal total aplicado incluyendo recirculación}

El caudal de recirculación requiere bombeo desde la salida del biofiltro (o del clarificador secundario) hacia la entrada del distribuidor rotatorio. La potencia requerida para este bombeo debe considerarse en el balance energético de la planta.

Para el dimensionamiento, se establece una carga orgánica volumétrica de diseño $C_v = 1.00$ kg DBO/m³·d, apropiada para biofiltros de alta carga con recirculación. El volumen de medio resulta:

\begin{equation}
V = \frac{W}{C_v} = \frac{28.78}{1.00} = 28.8\ \text{m}^3
\end{equation}
\captionequation{Volumen de medio filtrante por criterio de carga orgánica -- Ruta B}

El área superficial y el diámetro equivalente del biofiltro resultan:

\begin{equation}
A_s = \frac{V}{D} = \frac{28.8}{3.00} = 9.59\ \text{m}^2 \quad \Rightarrow \quad D_{\text{bf}} = \sqrt{\frac{4A_s}{\pi}} = 3.50\ \text{m}
\end{equation}

La carga hidráulica superficial total, considerando el caudal con recirculación, es:

\begin{equation}
\text{CHS} = \frac{Q_T}{A_s} = \frac{41.54}{9.59} = 4.330\ \text{m}^3\text{/m}^2\text{·h}
\end{equation}
\captionequation{Carga hidráulica superficial total con recirculación}

Este valor de CHS garantiza un riego intensivo del medio, manteniendo la biopelícula húmeda y activa incluso durante períodos de bajo caudal.

\subsubsection*{Geometría Completa del Biofiltro}

La altura total del biofiltro incluye componentes constructivas adicionales al medio filtrante, necesarias para el funcionamiento hidráulico y la operación del sistema. El desglose de alturas se presenta en la Tabla~\ref{tab:alturas_taf}.

\begin{table}[H]
\centering
\caption{Desglose de alturas constructivas -- Biofiltro TAF}
\label{tab:alturas_taf}
\begin{tabular}{lccc}
\toprule
\textbf{Componente} & \textbf{Símbolo} & \textbf{Altura (m)} & \textbf{Función} \\
\midrule
Espacio distribuidor--medio & $H_{\text{dist}}$ & 0.20 & Distribución hidráulica \\
Profundidad del medio & $D$ & 3.00 & Soporte biopelícula \\
Sistema underdrain & $H_{\text{under}}$ & 0.50 & Drenaje y ventilación \\
Bordo libre & $H_{\text{borde}}$ & 0.30 & Seguridad operacional \\
\midrule
\textbf{Altura total} & $\mathbf{H_{\text{total}}}$ & $\mathbf{4.00}$ & -- \\
\bottomrule
\end{tabular}
\end{table}

El espacio entre el distribuidor rotatorio y la superficie del medio ($H_{\text{dist}} = 0.20$ m) permite que el agua se distribuya en forma de abanico antes de contactar el medio, garantizando un mojado uniforme. El sistema underdrain ($H_{\text{under}} = 0.50$ m) recolecta el efluente tratado y permite la entrada de aire para ventilación natural. El bordo libre ($H_{\text{borde}} = 0.30$ m) proporciona seguridad contra salpicaduras y alberga el crecimiento de biopelícula en la superficie.

\subsubsection*{Sistema de Distribución Rotatorio}

El sistema de distribución rotatorio es el componente que aplica el agua residual uniformemente sobre la superficie del medio filtrante. Consiste en una columna central hueca (pivot) conectada a brazos radiales que giran por acción de la fuerza de reacción del agua al salir por las boquillas.

Para este diseño se utilizan 2 brazos distribuidores, seleccionados en función del diámetro del biofiltro. El caudal total que llega al distribuidor es $Q_{\text{aplicado}} = 41.54$ m³/h, resultando en un caudal por brazo de:

\begin{equation}
Q_{\text{brazo}} = \frac{Q_{\text{aplicado}}}{N_{\text{brazos}}} = \frac{41.54}{2} = 20.77\ \text{m}^3\text{/h} = 5.77\ \text{L/s}
\end{equation}
\captionequation{Caudal por brazo del distribuidor rotatorio}

La longitud de cada brazo corresponde al radio del biofiltro:

\begin{equation}
L_{\text{brazo}} = \frac{D_{\text{bf}}}{2} = \frac{3.50}{2} = 1.75\ \text{m}
\end{equation}

Cada brazo dispone de 8 boquillas distribuidas uniformemente, con un caudal por boquilla de 0.72 L/s. El diámetro de orificio calculado para una velocidad de salida de 2.0 m/s resulta $d_{\text{orificio}} = 21.4$ mm, valor dentro del rango constructivo estándar para este tipo de distribuidores.

La velocidad de rotación del distribuidor depende del caudal aplicado y del par de reacción generado por las boquillas. Velocidades típicas oscilan entre 0.5 y 2.0 rpm, garantizando una distribución uniforme sin generar salpicaduras excesivas.

\subsubsection{Verificación}

\subsubsection*{Verificación de Eficiencia de Remoción}

El dimensionamiento por carga orgánica volumétrica establece el tamaño del biofiltro a partir de criterios empíricos, pero no garantiza por sí solo que el sistema alcanzará la calidad de efluente requerida. Por esta razón, es obligatoria la verificación de eficiencia mediante el modelo cinético correspondiente (NRC para Ruta A, Germain para Ruta B), que incorpora explícitamente la temperatura, la profundidad del medio y la tasa hidráulica aplicada.

Para la Ruta B, aplicando el modelo de Germain con los parámetros del diseño, primero se corrige la constante cinética por temperatura:

\begin{equation}
k_T = k_{20} \cdot \theta^{(T-20)} = 0.068 \times 1.035^{(25.6-20)} = 0.0824\ \text{m/h}
\end{equation}
\captionequation{Corrección de constante cinética por temperatura}

La constante cinética corregida $k_T = 0.0824$ m/h refleja el aumento de la actividad microbiana a la temperatura de operación de 25.6°C, representando aproximadamente un 21.2\% de incremento respecto al valor de referencia a 20°C.

Aplicando el modelo de Germain:

\begin{equation}
S_e = S_0' \cdot \exp\!\left(-\frac{k_T \cdot D}{Q_T^n}\right) = 28.9 \times \exp\!\left(-\frac{0.0824 \times 3.00}{(4.330)^{0.5}}\right) = 25.6\ \text{mg/L}
\end{equation}
\captionequation{Concentración efluente estimada por modelo de Germain}

La eficiencia global del sistema, considerando la recirculación, resulta:

\begin{equation}
E = \frac{S_0 - S_e}{S_0} = \frac{57.7 - 25.6}{57.7} = 0.556 \approx 55.6\%
\end{equation}

La eficiencia calculada de 55.6\% refleja el desempeño esperado del biofiltro con los parámetros de diseño adoptados.

\textbf{Resultado de verificación:} La DBO₅ efluente calculada es 25.6 mg/L.

\subsubsection*{Verificación de Carga Orgánica Superficial}

La carga orgánica superficial (COS) representa la masa de DBO$_5$ aplicada por unidad de volumen del medio por día. Es un parámetro crítico porque determina la densidad de biomasa que puede mantenerse activa en el sistema y el riesgo de colmatación por acumulación excesiva de biopelícula.

\begin{equation}
\text{COS} = \frac{W}{V} = \frac{28.78}{28.8} = 1.000\ \text{kg DBO/m}^3\text{·d}
\end{equation}
\captionequation{Carga orgánica superficial del biofiltro}

El valor calculado de COS = 1.000 kg/m³·d se encuentra dentro del rango recomendado para biofiltros con medio plástico (hasta 1.20 kg/m³·d). Esta carga garantiza el desarrollo de una biopelícula activa sin riesgo de colmatación prematura del medio.

\subsubsection*{Verificación de Carga Hidráulica Superficial}

La carga hidráulica superficial (CHS) determina la frecuencia de contacto entre el agua residual y la biopelícula, así como la velocidad de percolación a través del medio. Valores bajos pueden causar desecación de la biopelícula durante períodos de bajo caudal, mientras que valores altos pueden producir arrastre de biomasa y canalización del flujo.

\begin{equation}
\text{CHS} = \frac{Q_{\text{aplicado}}}{A_s} = 4.330\ \text{m}^3\text{/m}^2\text{·h} = 103.92\ \text{m}^3\text{/m}^2\text{·d}
\end{equation}
\captionequation{Carga hidráulica superficial del biofiltro}

La CHS calculada de 4.330 m³/m²·h (103.92 m/d) se encuentra fuera del rango operacional recomendado (0.5--4.0 m³/m²·h). Se recomienda ajustar el área o la recirculación para alcanzar valores dentro del rango óptimo.

\subsubsection*{Tiempo de Retención Hidráulico}

El tiempo de retención hidráulico (TRH) representa el tiempo medio que el agua residual permanece en contacto con la biopelícula. Aunque en biofiltros el TRH es corto (típicamente menos de 1 hora), este parámetro es indicativo de la intensidad del contacto entre el sustrato y los microorganismos.

\begin{equation}
\text{TRH} = \frac{V}{Q_{\text{aplicado}}} = \frac{28.8}{41.54} = 0.69\ \text{h}
\end{equation}
\captionequation{Tiempo de retención hidráulico en el biofiltro}

El TRH de 0.69 horas es característico de biofiltros de medio plástico con las cargas hidráulicas y orgánicas de este proyecto. A pesar de ser relativamente corto, la alta densidad de biopelícula y la eficiente transferencia de masa en el medio estructurado permiten lograr las eficiencias de remoción verificadas.

\subsubsection*{Producción de Lodos Biológicos}

El biofiltro genera lodos biológicos (humus) por el desprendimiento periódico de la biopelícula excedente. Esta producción de sólidos debe considerarse en el diseño del sedimentador secundario y en el manejo de lodos de la planta.

La producción de humus puede estimarse como una fracción de la DBO$_5$ removida:

\begin{equation}
P_{\text{humus}} = Y \cdot \Delta\text{DBO} = 0.15 \times 16.00 = 2.40\ \text{kg SST/d}
\end{equation}
\captionequation{Producción de humus en el biofiltro TAF}

Donde $Y = 0.15$ kg SST/kg DBO es el coeficiente de producción de humus típico para biofiltros, y $\Delta\text{DBO} = 16.00$ kg/d es la DBO$_5$ removida en el sistema.

El lodo de humus producido en el biofiltro es relativamente estable aeróbicamente y puede ser manejado junto con los lodos del sedimentador secundario. La producción estimada de 2.40 kg SST/d representa una carga adicional de sólidos que debe acumularse o disponerse adecuadamente.

\subsubsection{Resultados}

La Tabla~\ref{tab:resumen_taf} consolida los parámetros geométricos, hidráulicos, de carga y de calidad del efluente del biofiltro TAF dimensionado, permitiendo una lectura integrada de los resultados del diseño.

\begingroup
\small
\begin{longtable}{ll}
\caption{Resumen de resultados del dimensionamiento -- Biofiltro TAF (Ruta B)}
\label{tab:resumen_taf}\\
\toprule
\textbf{Parámetro} & \textbf{Valor} \\
\midrule
\endfirsthead
\caption[]{Resumen de resultados del dimensionamiento -- Biofiltro TAF (Ruta B) (continuación)}\\
\toprule
\textbf{Parámetro} & \textbf{Valor} \\
\midrule
\endhead
\midrule
\multicolumn{2}{r}{\textit{Continúa en la siguiente página}} \\
\endfoot
\bottomrule
\endlastfoot
\multicolumn{2}{l}{\textit{Características del afluente}} \\
Caudal promedio & 498.5 m³/d \\
Caudal máximo horario & 1380.9 m³/d \\
DBO$_5$ afluente & 57.7 mg/L \\
Carga orgánica & 28.78 kg DBO/d \\
\midrule
\multicolumn{2}{l}{\textit{Metodología de diseño}} \\
Ruta seleccionada & B (Germain (con recirculación)) \\
Criterio de selección & COS_prelim = 0.60 kg/m³·d > 0.40 → Ruta B \\
\midrule
\multicolumn{2}{l}{\textit{Parámetros de recirculación (Ruta B)}} \\
Razón de recirculación & $R = 1.0$ \\
DBO afluente diluida & $S_0' = 28.9$ mg/L \\
Caudal de recirculación & $Q_R = 498.5$ m³/d \\
Caudal total aplicado & $Q_T = 41.54$ m³/h \\
\midrule
\multicolumn{2}{l}{\textit{Dimensionamiento}} \\
Volumen de medio & 28.8 m³ \\
Área superficial & 9.59 m² \\
Diámetro & 3.50 m \\
Profundidad del medio & 3.00 m \\
Altura total & 4.00 m \\
\midrule
\multicolumn{2}{l}{\textit{Cargas de operación}} \\
Carga orgánica volumétrica (COS) & 1.000 kg/m³·d \\
Carga hidráulica superficial (CHS) & 4.330 m³/m²·h \\
Tiempo de retención hidráulico & 0.69 h \\
\midrule
\multicolumn{2}{l}{\textit{Desempeño esperado}} \\
Eficiencia de remoción DBO$_5$ & 55.6\% \\
DBO$_5$ efluente calculada & 25.6 mg/L \\
DBO$_5$ removida & 16.00 kg/d \\
\midrule
\multicolumn{2}{l}{\textit{Sistema de distribución}} \\
Número de brazos & 2 \\
Caudal por brazo & 20.77 m³/h \\
Longitud de brazo & 1.75 m \\
Número de boquillas/brazo & 8 \\
Diámetro de orificio & 21.4 mm \\
\midrule
\multicolumn{2}{l}{\textit{Producción de lodos}} \\
Humus generado & 2.40 kg SST/d \\
SST efluente estimado & 48.5 mg/L \\
\end{longtable}
\endgroup

\subsubsection*{Conclusiones del Diseño}

El biofiltro de carga mecanizada hidráulica ha sido dimensionado siguiendo la \textbf{Ruta B} del manual técnico, utilizando el Germain (con recirculación). La geometría principal del sistema resulta:

\begin{itemize}[noitemsep,leftmargin=2em]
    \item Diámetro: 3.50 m
    \item Altura total: 4.00 m
    \item Volumen de medio filtrante: 28.8 m³
    \item Área superficial: 9.59 m²
\end{itemize}

El desempeño esperado del sistema, verificado mediante el modelo cinético correspondiente, indica:

\begin{itemize}[noitemsep,leftmargin=2em]
    \item Eficiencia de remoción de DBO$_5$: 55.6\%
    \item DBO$_5$ del efluente: 25.6 mg/L

\end{itemize}

La DBO₅ efluente calculada es 25.6 mg/L. El sistema incluye distribución rotatoria con 2 brazos, sistema underdrain para drenaje y ventilación natural, y produce aproximadamente 2.40 kg SST/d de humus que requiere separación en sedimentador secundario.


La figura siguiente presenta el corte transversal del biofiltro TAF, con sus componentes principales: distribuidor rotatorio, medio filtrante con biopelícula, sistema underdrain con ventilación natural y, conexiones hidráulicas y bomba de recirculación.

\begin{figure}[H]
\centering
\includegraphics[width=\textwidth]{figuras/Esquema_Biofiltro_TAF_RutaB.png}
\caption{Esquema del biofiltro TAF -- Ruta B -- Germain (con recirculación). Diámetro: 3.50~m, profundidad de medio: 3.00~m, altura total: 4.00~m, CHS: 4.330~m³/m²·h, eficiencia DBO$_5$: 56\%.}
\label{fig:esquema_biofiltro_taf_rutab}
\end{figure}

El esquema ilustra el recorrido del agua residual desde su entrada al distribuidor rotatorio, la percolación descendente a través del medio filtrante con biopelícula aerobia, y la recolección del efluente tratado en el sistema underdrain. Las aberturas laterales permiten la circulación natural de aire de abajo hacia arriba, manteniendo las condiciones aerobias necesarias para la biopelícula. La bomba de recirculación devuelve una fracción $R=1.0$ del efluente a la entrada del distribuidor, diluyendo el afluente y aumentando la carga hidráulica sobre el medio.

\subsection{Sedimentador Secundario}

\subsubsection{Dimensionamiento}

\textbf{Fundamentos teóricos y principios de operación}

El sedimentador secundario constituye una unidad de separación física por gravedad diseñada específicamente para remover los sólidos biológicos desprendidos de procesos de tratamiento aerobio previos. A diferencia de los sedimentadores primarios que remueven sólidos suspendidos en el agua cruda, el sedimentador secundario debe manejar flóculos biológicos activos (conocidos técnicamente como ``humus'' o lodos biológicos secundarios) que presentan características físicas y de sedimentación distintivas.

Los flóculos biológicos provenientes de filtros percoladores, MBBR o biofiltros aerados son agregados relativamente ligeros con densidades ligeramente superiores al agua, típicamente en el rango de 1,02--1,05 g/cm³, y velocidades de sedimentación que oscilan entre 0,3 y 1,2 m/h según Sperling \cite{sperling2007}. Estas partículas son susceptibles a la compactación y pueden re-suspenderse fácilmente si la velocidad ascensional del líquido excede ciertos límites críticos, lo que justifica el diseño conservador de estos tanques.

\textbf{Densidad de flóculos y su relevancia en el diseño:}

La densidad de los flóculos biológicos es un parámetro crítico que determina la velocidad de sedimentación y, consecuentemente, la eficiencia de separación en el sedimentador secundario. A diferencia de los sólidos minerales (arena, grava) que tienen densidades superiores a 2,5 g/cm³, los flóculos biológicos presentan:

\begin{itemize}[noitemsep,leftmargin=2em]
    \item \textbf{Densidad aparente:} 1,02--1,05 g/cm³ (incluyendo agua intraflocular), solo ligeramente superior al agua (1,0 g/cm³).
    \item \textbf{Densidad real de sólidos:} 1,10--1,15 g/cm³ si se considera únicamente la materia seca de los microorganismos.
    \item \textbf{Porosidad del flóculo:} 90--95\% de agua embebida en la estructura del biofloc.
\end{itemize}

Esta baja diferencia de densidad (2--5\% respecto al agua) implica que la fuerza de flotación es casi igual al peso del flóculo, resultando en velocidades de sedimentación muy bajas. Por esta razón, el diseño del sedimentador secundario requiere:

\begin{itemize}[noitemsep,leftmargin=2em]
    \item \textbf{SOR conservadora:} Valores de 16--32 m³/m²·d (equivalente a 0,67--1,33 m/h) que garanticen que la velocidad ascensional del líquido no supere la velocidad de caída de los flóculos.
    \item \textbf{Tiempo de retención suficiente:} Mínimo 1,5 horas para permitir que incluso los flóculos más pequeños sedimenten.
    \item \textbf{Gradientes de velocidad bajos:} Evitar turbulencias que puedan fragmentar los flóculos o mantenerlos en suspensión.
\end{itemize}

\textit{Nota técnica:} La densidad de los flóculos puede variar según la carga orgánica aplicada al sistema biológico previo. A mayor carga, los flóculos tienden a ser más livianos (mayor contenido de biomasa activa, menor densidad), requiriendo diseños más conservadores del sedimentador.

\textbf{Mecanismos de sedimentación y transporte}

El proceso de separación en el sedimentador secundario involucra cuatro mecanismos principales: (1) \textit{sedimentación discreta} de partículas individuales en la zona superior, (2) \textit{sedimentación en floculación} donde las partículas se agregan durante su descenso aumentando su velocidad de caída, (3) \textit{compresión por estratificación} en la zona de acumulación de lodos, y (4) \textit{transporte por flujo laminar} hacia el sistema de recolección de efluente clarificado.

El diseño geométrico del tanque debe asegurar un flujo hidrodinámico que minimice la turbulencia y los cortocircuitos hidráulicos. La relación diámetro/profundidad típica para sedimentadores circulares secundarios varía entre 1,5:1 y 4:1, según recomendaciones de Metcalf \& Eddy \cite{metcalf2014}. Un diseño más profundo favorece el almacenamiento de lodos y reduce el riesgo de re-suspensión por venteo, mientras que un diámetro mayor mejora la eficiencia de sedimentación al reducir la velocidad ascensional superficial.

\textbf{Características del humus y requerimientos de almacenamiento}

El humus o lodo biológico secundario acumulado en el fondo del sedimentador presenta una concentración típica de sólidos totales entre 2\% y 4\% en peso, con un índice de volumen de lodos (IVL) que puede variar significativamente según la carga orgánica aplicada al sistema biológico precedente. Según Sperling \cite{sperling2007}, el lodo de filtros percoladores presenta mejores características de sedimentación que el de lodos activados, con valores de IVL típicamente menores a 150 mL/g, lo que permite diseños más compactos.

La capacidad de almacenamiento de lodos debe dimensionarse considerando: (a) la producción diaria de humus estimada en 0,2--0,5 kg SST por kg de DBO₅ removida en el sistema biológico, (b) la frecuencia de purga programada (típicamente diaria o semanal según operación), y (c) un factor de seguridad del 25\% para manejar picos de producción de sólidos.

\textbf{Criterios de diseño y parámetros fundamentales}

El diseño del sedimentador secundario se fundamenta en tres criterios interrelacionados: (1) la \textit{tasa de desbordamiento superficial} (SOR) como parámetro de control de la sedimentación, (2) el \textit{tiempo de retención hidráulico} (TRH) como garantía de tiempo de contacto suficiente para la separación, y (3) la \textit{carga sobre el vertedero perimetral} como control del flujo de salida que evite arrastre de sólidos.

La SOR representa la velocidad ascensional teórica del líquido en el tanque y se calcula como el cociente entre el caudal y el área superficial neta. Para sedimentadores secundarios posteriores a procesos de biopelícula, Metcalf \& Eddy \cite{metcalf2014} establecen un rango operativo óptimo de 16--32 m³/m²·d para condiciones de caudal medio, con un límite máximo absoluto de 45 m³/m²·d durante picos horarios para evitar la re-suspensión del lodo sedimentado.

El TRH mínimo requerido es de 1,5 horas para garantizar que los flóculos biológicos tengan tiempo suficiente para sedimentar antes de ser arrastrados hacia el efluente. Valores superiores a 4 horas son aceptables y favorecen la clarificación, especialmente en climas tropicales donde las temperaturas elevadas pueden afectar la densidad del agua y los flóculos.

La carga sobre vertedero perimetral debe mantenerse por debajo de 250 m³/m·d para sedimentadores circulares secundarios, aunque valores más conservadores de 186 m³/m·d son recomendados por Sperling \cite{sperling2007} para sistemas de biopelícula, garantizando una velocidad de aproximación al vertedero suficientemente baja que no arrastre los sólidos sedimentados hacia el efluente clarificado.

\textbf{Parámetros de diseño}

\begin{table}[H]
\centering
\caption{Parámetros de diseño del sedimentador secundario}
\footnotesize
\setlength{\tabcolsep}{3pt}
\renewcommand{\arraystretch}{1.18}
\begin{tabular}{@{}p{0.25\textwidth}p{0.17\textwidth}p{0.25\textwidth}p{0.17\textwidth}@{}}
\toprule
\textbf{Parámetro} & \textbf{Valor adoptado} & \textbf{Rango/Unidad} & \textbf{Fuente} \\
\midrule
Caudal medio por línea (Q) & 498.5 m³/d & – & – \\
Factor de pico horario ($f_p$) & 2.8 & – & Metcalf \& Eddy \\
Caudal máximo horario ($Q_{max}$) & 1380.9 m³/d & – & – \\
Tasa de desbordamiento\newline superficial (SOR) & 16.0 m³/m²·d & 16--32 m³/m²·d & Metcalf \& Eddy \\\\
Profundidad lateral & 3.50 m & 3.0--4.5 m & Sperling \\
Tiempo de retención\newline hidráulico (HRT) & 5.2 h & $\geq$ 1.5 h & Metcalf \& Eddy \\\\
SOR máximo\newline permisible & Cumple & $\leq$ 45 m³/m²·d & Metcalf \& Eddy \\\\
Carga sobre vertedero\newline perimetral & 25.2 m³/m·d & $\leq$ 250 m³/m·d & Metcalf \& Eddy \\\\
Zona de almacenamiento\newline de lodos & 0.50 m (tolva) & – & Metcalf \& Eddy \\\\
Bordo libre & 0.30 m & – & Norma constructiva \\
\bottomrule
\end{tabular}
\end{table}

El diseño del sedimentador secundario se fundamenta en la tasa de desbordamiento superficial (SOR), que relaciona el caudal con el área superficial del tanque. Para sedimentadores secundarios ubicados después de procesos biológicos aerobios (como filtros percoladores, MBBR o biofiltros), Metcalf \& Eddy (2014) recomiendan valores de SOR entre 16 y 32 m³/m²·d para operación normal, con un límite máximo de 45 m³/m²·d para condiciones de pico horario.

El tiempo de retención hidráulico (TRH) resultante del volumen del tanque debe ser al menos 1.5 horas para garantizar la sedimentación efectiva de los sólidos biológicos. El valor calculado para este diseño es 5.2 horas.

\textbf{Carga Superficial -- Dimensionamiento}

Según Metcalf \& Eddy \cite{metcalf2014}, el dimensionamiento de sedimentadores secundarios se fundamenta en la tasa de desbordamiento superficial (SOR), que representa la relación entre el caudal y el área superficial disponible para la sedimentación. Este parámetro garantiza que la velocidad ascensional del líquido sea suficientemente baja para permitir la decantación gravitacional de los flóculos biológicos. La ecuación fundamental es:

\begin{equation}
A_s = \frac{Q}{SOR}
\end{equation}
\captionequation{Area superficial del sedimentador - Criterio SOR}

\textit{Donde:}
\begin{itemize}[noitemsep,leftmargin=2em]
    \item[$A_s$] = Área superficial del sedimentador (m²)
    \item[$Q$] = Caudal diario (498.5 m³/d)
    \item[$SOR$] = Tasa de desbordamiento superficial (16.0 m³/m²·d)
\end{itemize}

\begin{equation}
A_s = \frac{498.5}{16.0} = 31.16 \text{ m}^2
\end{equation}

Correspondiendo a un diámetro de 6.30 m para configuración circular. La profundidad lateral de 3.50 m proporciona un volumen útil de 109.1 m³.

\textbf{Desglose de altura total:}

\begin{table}[H]
\centering
\small
\begin{tabular}{lc}
\toprule
Zona & Altura \\
\midrule
Zona de sedimentación (profundidad lateral) & 3.50 m \\
Zona de almacenamiento de lodos (tolva) & 0.50 m \\
Bordo libre & 0.30 m \\
\midrule
\textbf{Altura total de construcción} & \textbf{4.30 m} \\
\bottomrule
\end{tabular}
\end{table}

\textbf{Vertedero Perimetral -- Dimensionamiento}

El vertedero perimetral constituye el elemento de recolección del efluente clarificado. Su diseño debe garantizar una distribución uniforme del caudal perimetral para evitar zonas de acumulación o cortocircuito. Según Metcalf \& Eddy \cite{metcalf2014} y Sperling \cite{sperling2007}, la carga sobre el vertedero (weir loading) es un parámetro crítico que determina la uniformidad del flujo de salida. Para un sedimentador circular, el perímetro se calcula como:

\begin{equation}
P = \pi \times D = \pi \times 6.30 = 19.79 \text{ m}
\end{equation}

La carga sobre el vertedero se calcula como:

\begin{equation}
q_{\text{vertedero}} = \frac{Q}{P} = \frac{498.5}{19.79} = 25.2 \text{ m}^3\text{/m·d}
\end{equation}

\textbf{Almacenamiento de Lodos -- Dimensionamiento}

La producción de humus —sólidos biológicos desprendidos de la etapa biológica aerobia como resultado del control de espesor de la biopelícula— se estima siguiendo los criterios de Metcalf \& Eddy \cite{metcalf2014}. El factor de producción de 0.15 kg SST por kg de DBO removida representa la fracción de sólidos generados en el proceso de oxidación biológica que deben ser separados en el sedimentador secundario. Este valor es característico de sistemas con biopelícula aerobia y permite estimar la carga de sólidos sobre el área de sedimentación:

\begin{equation}
P_{\text{humus}} = 0.15 \times \text{DBO}_{\text{removida}} = 2.4 \text{ kg SST/d}
\end{equation}

La tasa de aplicación de sólidos sobre el área del sedimentador resulta:

\begin{equation}
C_{\text{sólidos}} = \frac{P_{\text{humus}}}{A_s} = \frac{2.4}{31.16} = 0.08 \text{ kg SST/m}^2\text{·d}
\end{equation}

\subsubsection{Verificación}

\textbf{Verificación de Carga Superficial}

De acuerdo con los criterios de Metcalf \& Eddy \cite{metcalf2014} y WEF MOP-8 \cite{wef_mop8_2010}, el sedimentador debe verificarse para condiciones de caudal máximo horario. El factor de pico ($f_p$) de 2.8 representa la relación entre el caudal máximo horario y el caudal medio diario, típico de sistemas de alcantarillado municipal. La verificación establece que:

\begin{equation}
Q_{\max} = f_p \times Q = 2.8 \times 498.5 = 1380.9 \text{ m}^3\text{/d}
\end{equation}

A caudal máximo, la tasa de desbordamiento resulta:

\begin{equation}
SOR_{\max} = \frac{Q_{\max}}{A_s} = \frac{1380.9}{31.16} = 44.3 \text{ m}^3\text{/m}^2\text{·d}
\end{equation}

la tasa de desbordamiento máxima (44.3 m³/m²·d) está dentro del rango operativo seguro ($\leq$ 45 m³/m²·d) con un margen del 1.5%. El margen de seguridad respecto al límite de 45 m³/m²·d es del 1.5 por ciento. Según Metcalf & Eddy (2014), mantener el SOR máximo por debajo del límite asegura que la sedimentación no se vea afectada por las condiciones de pico, preservando la calidad del efluente.

\textbf{Verificación del Vertedero Perimetral}

La carga calculada de 25.2 m³/m·d se compara contra el límite establecido por Metcalf & Eddy de 250 m³/m·d para sedimentadores secundarios. Este límite surge de la necesidad de mantener velocidades de aproximación al vertedero suficientemente bajas que no arrastren los sólidos ya sedimentados hacia el efluente. Con la carga calculada, el estado de la verificación es Cumple.

\textbf{Verificación de Almacenamiento de Lodos}

La tasa de aplicación de sólidos calculada de 0.08 kg SST/m²·d se compara contra el límite conservador de 100 kg SST/m²·d adoptado para este diseño, basado en criterios de Metcalf & Eddy para sedimentadores secundarios después de procesos biológicos aerobios. El rango operativo típico de referencia es 50--150 kg SST/m²·d. El estado de la verificación es Cumple.

Respecto al tiempo de retención a caudal mínimo (factor 0.4 del medio), el TRH resultante es de 13.1 horas. Este valor excede 8.0 horas, por lo que se recomienda monitorear la operación para evitar condiciones sépticas.

\textbf{Nota técnica sobre operación a caudal mínimo:} El HRT de 13.1 h a caudal mínimo puede favorecer condiciones sépticas en el fondo del sedimentador si la operación no se controla. Se recomiendan las siguientes medidas operativas: (1) mantener una recirculación interna del efluente clarificado hacia la entrada del sedimentador para mantener la turbulencia y prevenir la putrefacción de lodos; (2) instalar un sistema de control de nivel que permita operar con tirante variable, reduciendo el volumen almacenado en horas de bajo caudal; (3) programar purgas periódicas de lodos acumulados en la tolva, mínimo cada 48 horas. Estas medidas deberán detallarse en la ingeniería de detalle del proyecto.

\subsubsection{Resultados}

\begingroup
\small
\begin{longtable}{ll}
\caption{Dimensiones del sedimentador secundario}\\
\toprule
Parámetro & Valor \\
\midrule
\endfirsthead
\caption[]{Dimensiones del sedimentador secundario (continuación)}\\
\toprule
Parámetro & Valor \\
\midrule
\endhead
\midrule
\multicolumn{2}{r}{\textit{Continúa en la siguiente página}} \\
\endfoot
\bottomrule
\endlastfoot
Diámetro & 6.30 m \\
Área superficial & 31.16 m² \\
Profundidad lateral & 3.50 m \\
Altura total de construcción & 4.30 m \\
Volumen útil & 109.1 m³ \\
\midrule
\multicolumn{2}{l}{\textit{Parámetros hidráulicos}} \\
\midrule
SOR (promedio) & 16.0 m³/m²·d \\
SOR (caudal máximo) & 44.3 m³/m²·d \\
HRT (promedio) & 5.2 h \\
HRT (caudal máximo) & 1.9 h \\
Carga sobre vertedero & 25.2 m³/m·d \\
Carga de sólidos & 0.08 kg/m²·d \\
\end{longtable}
\endgroup

\begin{figure}[H]
\centering
\includegraphics[width=0.92\textwidth]{figuras/Esquema_Sedimentador_Secundario.png}
\caption{Esquema del sedimentador secundario circular: entrada central del efluente de la etapa biológica, zona de clarificación, tolva de lodos, vertedero perimetral y salida de efluente clarificado. Diámetro adoptado: 6.30 m, profundidad lateral: 3.50 m, SOR promedio: 16.0 m³/m²·d.}
\label{fig:sedimentador_secundario}
\end{figure}

Se debe incluir un sistema de recolección de lodos (rastrillos o succionadores) y un mecanismo de extracción de lodos desde el fondo. Se recomienda considerar el caudal mínimo en el diseño operacional para evitar estancamiento.

\subsection{Sistema de Desinfeccion con Hipoclorito de Sodio}





La desinfección constituye la etapa final del proceso de tratamiento de aguas residuales, cuyo objetivo es la inactivación de microorganismos patógenos (bacterias, virus y protozoos) presentes en el efluente tratado, garantizando que su descarga al cuerpo receptor no represente riesgo para la salud pública ni para los ecosistemas acuáticos. En el contexto de la normativa ecuatoriana, la TULSMA (Texto Unificado de Legislación Secundaria del Ministerio del Ambiente) establece un límite máximo de 3000 NMP/100mL para coliformes fecales en vertimientos de aguas residuales tratadas.

El mecanismo de desinfección con compuestos clorados se fundamenta en la oxidación de componentes celulares esenciales de los microorganismos. El hipoclorito de sodio (NaOCl), al disolverse en agua, establece un equilibrio químico que genera ácido hipocloroso (HOCl) según la siguiente reacción:

\begin{equation}
\text{NaOCl} + \text{H}_2\text{O} \rightleftharpoons \text{HOCl} + \text{Na}^+ + \text{OH}^-
\end{equation}
\captionequation{Disolución del hipoclorito de sodio en agua}

El ácido hipocloroso (HOCl) es la especie biocida predominante y efectiva, siendo aproximadamente 80--100 veces más potente como desinfectante que el ion hipoclorito (OCl$^-$). La especiación entre HOCl y OCl$^-$ depende críticamente del pH del agua, según el siguiente equilibrio:

\begin{equation}
\text{HOCl} \rightleftharpoons \text{H}^+ + \text{OCl}^- \quad (\text{pKa} \approx 7,5 \text{ a 25°C})
\end{equation}
\captionequation{Equilibrio ácido-base del ácido hipocloroso}

A pH inferior a 7,5 predomina el HOCl (forma más efectiva), mientras que a pH superior a 7,5 predomina el OCl$^-$. Para el rango típico de pH de aguas residuales tratadas (6,5--8,5), se asegura una fracción significativa de HOCl que garantiza la eficacia desinfectante. La temperatura del agua también afecta este equilibrio: a mayor temperatura, el valor de pKa disminuye ligeramente, favoreciendo la forma HOCl.

\textbf{Mecanismos de inactivación microbiana}

La acción biocida del cloro se ejerce mediante múltiples mecanismos simultáneos: (1) \textit{oxidación de la pared celular}, alterando la permeabilidad y causando lisis celular; (2) \textit{denaturación de enzimas}, particularmente aquellas con grupos sulfhidrilo (-SH) que son oxidados a disulfuros; (3) \textit{daño al material genético}, mediante reacción con bases nitrogenadas del ADN y ARN; y (4) \textit{alteración del metabolismo celular}, inhibiendo la síntesis de ATP y otras funciones energéticas.

La susceptibilidad de los microorganismos al cloro varía significativamente. Las bacterias (incluyendo coliformes fecales) son generalmente las más susceptibles, con valores CT de 0,1--5 mg$\cdot$min/L para 2-log de reducción. Los virus requieren CT mayores (5--20 mg$\cdot$min/L), mientras que los protozoos como \textit{Giardia} y \textit{Cryptosporidium} son los más resistentes, requiriendo CT de 30--100 mg$\cdot$min/L o alternativas como UV o ozono.

\textbf{El concepto CT como parámetro de diseño}

El producto concentración-tiempo (CT) constituye el parámetro fundamental para el dimensionamiento de sistemas de desinfección. Fue desarrollado por la USEPA como criterio regulatorio y ha sido adoptado internacionalmente. Físicamente, el CT representa la exposición acumulada de los microorganismos al desinfectante, integrando tanto la concentración residual como el tiempo de contacto efectivo.

El CT requerido depende de: (a) el microorganismo objetivo y la reducción logarítmica deseada, (b) la temperatura del agua (a mayor temperatura, menor CT requerido debido a cinética más rápida), (c) el pH (afecta la especiación del cloro), y (d) la presencia de materia orgánica y otros consumidores de cloro (demanda).

Según Metcalf \& Eddy \cite{metcalf2014} y la USEPA \cite{usepa2003}, los valores CT típicos para inactivación de coliformes en efluentes secundarios son:

\begin{itemize}[noitemsep,leftmargin=2em]
    \item CT = 5--10 mg$\cdot$min/L: 2--3 log de reducción (99--99,9\%)
    \item CT = 10--20 mg$\cdot$min/L: 3--4 log de reducción (99,9--99,99\%)
    \item CT = 20--40 mg$\cdot$min/L: 4--5 log de reducción (99,99--99,999\%)
\end{itemize}

\textbf{Criterio de disenio basado en CF objetivo}

El dimensionamiento del sistema de desinfección se fundamenta en un criterio de diseño explícito basado en la concentración de coliformes fecales (CF) objetivo. A diferencia de un enfoque de dosis fija, este método garantiza que el sistema sea capaz de alcanzar el nivel de inactivación microbiológica requerido, partiendo de la carga de coliformes de entrada y definiendo la reducción logarítmica necesaria para cumplir con el objetivo de calidad del efluente.

Para este diseño, se establece un CF objetivo de 3000 NMP/100mL, correspondiente al límite establecido por la TULSMA para vertimiento de aguas residuales tratadas. Con una concentración de entrada de 1287 NMP/100mL, la reducción logarítmica requerida es:

\begin{equation}
\text{Log reducción}_{\text{req}} = \log_{10}\left(\frac{CF_{\text{entrada}}}{CF_{\text{objetivo}}}\right) = \log_{10}\left(\frac{1287}{3000}\right) = 0.00 \text{ log}
\end{equation}
\captionequation{Log reduccion requerida para cumplir CF objetivo}

El CT requerido para lograr esta reducción se calcula mediante el coeficiente de inactivación adoptado:

\begin{equation}
CT_{\text{req}} = \frac{\text{Log reducción}_{\text{req}}}{k} = \frac{0.00}{0.22} = 0.0 \text{ mg$\cdot$min/L}
\end{equation}
\captionequation{CT requerido para alcanzar la reduccion logaritmica objetivo}

\textbf{Filosofía de diseño adoptada: TRH base + Residual calculado}

El diseño adopta una filosofía de tiempo de contacto base y residual calculado. Primero se dimensiona el tanque de contacto con un TRH base de referencia (30 min), obteniéndose un TRH real de 53.7 min según las dimensiones constructivas adoptadas. Posteriormente, se calcula el residual requerido para alcanzar el CT objetivo con ese TRH real:

\begin{equation}
C_{\text{req}} = \frac{CT_{\text{req}}}{TRH_{\text{real}}} = \frac{0.0}{53.7} = 0.00 \text{ mg/L}
\end{equation}
\captionequation{Residual requerido para alcanzar CT objetivo con TRH real}

El residual adoptado corresponde al valor requerido calculado, optimizando la dosis para cumplir justo con el objetivo de desinfección:

\begin{equation}
C_{\text{adoptado}} = C_{\text{req}} = 0.56 \text{ mg/L}
\end{equation}
\captionequation{Residual adoptado para cumplir CF objetivo}

\textit{Nota operativa:} El residual calculado (0.00 mg/L) es inferior al mínimo operativo recomendado (0.5 mg/L) para garantizar margen de seguridad en la operación. Se recomienda monitorear periódicamente para asegurar que el residual efectivo se mantenga en el rango de 0,5--1,0 mg/L.

El CT real alcanzado con el residual adoptado y el TRH real es:

\begin{equation}
CT_{\text{real}} = C_{\text{adoptado}} \times TRH_{\text{real}} = 0.56 \times 53.7 = 30.0 \text{ mg$\cdot$min/L}
\end{equation}
\captionequation{CT real alcanzado}

Como criterio adicional de verificación, se establece un CT mínimo de 30 mg$\cdot$min/L, valor conservador que asegura el cumplimiento del límite de coliformes fecales establecido por la TULSMA, considerando las variaciones operativas esperadas.

\textbf{Demanda de cloro y dosificación}

La demanda de cloro representa la cantidad de cloro consumida por reacciones con sustancias presentes en el agua antes de que quede disponible cloro residual libre para desinfección. Los principales consumidores son: amoníaco (formando cloraminas), materia orgánica reducida, compuestos férricos y manganesos, y sulfuros. En efluentes secundarios tratados (post-filtro percolador), la demanda típica adoptada para referencia de diseño varía entre 2--5 mg/L dependiendo de la calidad del efluente.

El cloro residual es la concentración que permanece después de satisfacer la demanda y es la que ejerce el efecto desinfectante. El diseño debe garantizar un residual mínimo de 0.6 mg/L al final del tiempo de contacto para asegurar eficacia. La dosis total a aplicar es:

\begin{equation}
\text{Dosis total} = \text{Demanda} + \text{Residual objetivo}
\end{equation}
\captionequation{Descomposición de la dosis de cloro}

\textbf{Diseño del tanque de contacto}

El tanque de contacto se dimensiona para proporcionar un tiempo de retención hidráulico (TRH) adecuado que permita alcanzar el CT objetivo con el residual calculado. El TRH efectivo es el tiempo que permanece el agua en el tanque, calculado como el cociente entre el volumen útil y el caudal. Según Metcalf \& Eddy, para desinfección de efluentes secundarios se recomienda un TRH de 15--45 min; para este caso, se parte de un TRH base de 30 min para el dimensionamiento geométrico.

La geometría del tanque debe favorecer el flujo en pistón (plug flow) para evitar cortocircuitos hidráulicos que reduzcan el tiempo efectivo de contacto. Por ello, el volumen de contacto se organiza como una cámara con bafles alternados tipo culebrín, de modo que el recorrido hidráulico aumente y la relación recorrido/ancho del canal sea verificable.

\textbf{Parámetros de diseño adoptados}

El sistema se dimensiona considerando los siguientes parámetros fundamentales, establecidos según criterios bibliográficos y normativos:

\begin{table}[H]
\centering
\caption{Parámetros de diseño de la desinfección}
\small
\begin{tabular}{@{}p{0.28\textwidth}p{0.18\textwidth}p{0.22\textwidth}p{0.18\textwidth}@{}}
\toprule
\textbf{Parámetro} & \textbf{Valor adoptado} & \textbf{Rango/Unidad} & \textbf{Fuente} \\
\midrule
Demanda de cloro & 3.5 mg/L & 2--5 mg/L & Estimado efluente post-sedimentador \\
Residual requerido & 0.00 mg/L & -- & Cálculo desde CT \\
Residual adoptado & 0.56 mg/L & $>$0.00 mg/L & Para cumplir CF objetivo \\
Dosis total & 4.1 mg/L & 3--10 mg/L & Metcalf \& Eddy \cite{metcalf2014} \\
TRH base de diseño & 30 min & 15--45 min & Configuración \\
TRH real del tanque & 53.7 min & $\geq$ 30 min & Dimensionado \\
Relación recorrido/ancho & 30.0:1 & $\geq$ 30:1 & Criterio de flujo pistón \\
CT requerido & 0.0 mg$\cdot$min/L & -- & Desde log reducción \\
CT real & 30.0 mg$\cdot$min/L & $\geq$ 30 mg$\cdot$min/L & Residual $\times$ TRH \\
\bottomrule
\end{tabular}
\end{table}

La dosis total de cloro se descompone en la demanda del efluente (consumida por amoníaco y materia orgánica) más el residual requerido para desinfección:

\begin{equation}
\text{Dosis total} = \text{Demanda} + \text{Residual} = 3.5 + 0.6 = 4.1 \text{ mg/L}
\end{equation}
\captionequation{Dosis total de cloro}

\textbf{Cálculo del producto CT}

La eficacia de la desinfección se cuantifica mediante el parámetro CT:

\begin{equation}
CT = C \times t
\end{equation}
\captionequation{Producto concentracion-tiempo CT}

\textit{Donde:}
\begin{itemize}[noitemsep,leftmargin=2em]
    \item[$CT$] = Producto concentración-tiempo (mg$\cdot$min/L)
    \item[$C$] = Cloro residual (0.6 mg/L)
    \item[$t$] = Tiempo de contacto real en la cámara (53.7 min)
\end{itemize}

\begin{equation}
CT = 0.56 \times 53.7 = 30.0 \text{ mg$\cdot$min/L}
\end{equation}

La reducción de coliformes se estima mediante una relación empírica simplificada, válida como aproximación preliminar de diseño para efluentes secundarios típicos:

\begin{equation}
\text{Log reducción} \approx 0.22 \times CT
\end{equation}
\captionequation{Estimacion empirica de log-reduccion (simplificacion para diseño preliminar)}

\textit{Nota metodológica:} Esta relación lineal es una simplificación empírica calibrada para condiciones típicas de aguas residuales tratadas (pH 6.5--8.5, temperatura 20--25°C, baja materia orgánica residual). El modelo cinético formal de inactivación microbiana es el de Chick-Watson: $N/N_0 = e^{-k \cdot C^n \cdot t}$, donde $k$ es la constante de inactivación específica para cada microorganismo y $n$ es el coeficiente de dilución (típicamente $n \approx 1$ para cloro libre). La forma lineal utilizada es práctica para diseño preliminar pero debe verificarse con ensayos de desinfección específicos si se requiere precisión rigurosa.

\begin{equation}
\text{Log reducción} \approx 0.22 \times 30 = 6.6 \text{ log}
\end{equation}

Los coliformes finales se calculan como:

\begin{equation}
CF_{final} = \frac{CF_{inicial}}{10^{\text{Log reduccion}}} = \frac{1287}{10^{6.6}} = 0 \text{ NMP/100mL}
\end{equation}

\textbf{Dimensionamiento del tanque de contacto}

El volumen del tanque se determina por:

\begin{equation}
V = Q \times t
\end{equation}
\captionequation{Volumen del tanque de contacto}

\textit{Donde:}
\begin{itemize}[noitemsep,leftmargin=2em]
    \item[$V$] = Volumen del tanque (m³)
    \item[$Q$] = Caudal (498.5 m³/d = 0.346 m³/min)
    \item[$t$] = Tiempo de contacto base (30 min)
\end{itemize}

\begin{equation}
V_{min} = 0.346 \times 30 = 10.4 \text{ m}^3
\end{equation}

El volumen minimo por tiempo de contacto es 10.4 m3; la geometria serpentina adoptada entrega 18.6 m3 y un tiempo hidraulico teorico de 53.7 min.

La cámara se resuelve con 5 canales en culebrín, ancho hidráulico de 0.50 m por canal, bafles de 0.15 m y un recorrido hidráulico total de 15.0 m. Las dimensiones en planta resultantes son 3.0 m de largo por 3.1 m de ancho.

\textbf{Consumo de cloro y requerimientos de producto comercial}

El consumo de cloro activo (como Cl$_2$) se calcula como:

\begin{equation}
\text{Consumo Cl}_2 = \frac{D \times Q}{1000} = \frac{4.1 \times 498.5}{1000} = 2.02 \text{ kg Cl}_2\text{/d}
\end{equation}

\textbf{Conversión a hipoclorito de sodio comercial (NaOCl):}

El hipoclorito de sodio se comercializa típicamente al 10--12.5\% de cloro disponible. La conversión se realiza mediante:

\begin{equation}
\text{Consumo NaOCl} = \frac{\text{Consumo Cl}_2}{[\% \text{ NaOCl}]} = \frac{2.02}{10\%} = 20.2 \text{ kg NaOCl/d}
\end{equation}
\captionequation{Conversion a hipoclorito de sodio comercial}

Considerando una densidad de 1.10 kg/L para la solución al 10\%:

\begin{equation}
\text{Volumen NaOCl} = \frac{20.2 \text{ kg/d}}{1.10 \text{ kg/L}} = 18.4 \text{ L/d} \approx 552 \text{ L/mes}
\end{equation}

El volumen de almacenamiento en bodega se recomienda para 30 días de operación: 552 L (equivalente a 0.6 m³).

\textbf{Nota sobre subproductos de la desinfección}

La cloración de aguas residuales puede generar subproductos de desinfección (DBPs, por sus siglas en inglés), principalmente trihalometanos (THMs) y ácidos haloacéticos (HAAs), como resultado de la reacción del cloro con la materia orgánica residual presente en el efluente tratado.

Para plantas de tratamiento de pequeña escala como la proyectada, la formación de THMs puede controlarse mediante cuatro medidas operativas integrales. En primer lugar, el control de dosis debe enfocarse en aplicar la dosis mínima efectiva que garantice el residual requerido, evitando sobredosis innecesaria que incrementaría la generación de subproductos. En segundo lugar, el monitoreo de residual combinado debe mantenerse en el rango de 0.6--1.0 mg/L, ya que valores superiores favorecen la formación de DBPs. En tercer lugar, la calidad del efluente previo resulta determinante: asegurar que el tratamiento biológico previo (HAFV) remueva adecuadamente la materia orgánica reduce el potencial de formación de THMs, dado que una menor DQO/DBO residual implica menor disponibilidad de precursores orgánicos. Finalmente, el monitoreo periódico mediante análisis trimestral de THMs totales en el efluente final debe realizarse si la autoridad ambiental lo requiere; para vertimiento a cuerpos receptores no destinados a consumo humano, los requisitos normativos son menos estrictos que para agua potable.

Para el presente diseño, el riesgo de formación de THMs se considera manejable mediante la operación adecuada del sistema de dosificación y el control de residual, dado que el efluente del HAFV tiene baja concentración de materia orgánica precursora y la dosis de cloro aplicada es moderada.

\subsubsection{Verificación}

\textbf{Verificación de cumplimiento normativo}

El sistema de desinfección se verifica contra los criterios establecidos en la TULSMA y las recomendaciones de Metcalf \& Eddy para asegurar la eficacia del proceso.


El Texto Unificado de Legislación Secundaria del Ministerio del Ambiente (TULSMA) establece diferentes límites de coliformes fecales según el uso receptor del agua. La siguiente tabla compara el valor calculado de coliformes finales (0 NMP/100mL) con los límites establecidos para cada categoría de uso:

\begin{table}[H]
\centering
\caption{Comparación con límites TULSMA por uso del agua}
\small
\begin{tabular}{@{}llcc@{}}
\toprule
\textbf{Uso / Tabla TULSMA} & \textbf{CF (límite)} & \textbf{CF calculado} & \textbf{Dictamen} \\
\midrule
Agua dulce -- Tabla 12 & $\leq$3000 & 0 & \textbf{CUMPLE} \\
Agua marina -- Tabla 13 & $\leq$3000 & 0 & \textbf{CUMPLE} \\
Alcantarillado -- Tabla 11 & --- & 0 & \textbf{N/A} \\
Flora y fauna -- Tabla 3 & $\leq$200 & 0 & \textbf{CUMPLE} \\
Agrícola/riego -- Tabla 6 & $\leq$1000 & 0 & \textbf{CUMPLE} \\
Pecuario -- Tabla 7 & $\leq$1000 & 0 & \textbf{CUMPLE} \\
Recreativo primario -- Tabla 9 & $\leq$200 & 0 & \textbf{CUMPLE} \\
Recreativo secundario -- Tabla 10 & $\leq$2000 & 0 & \textbf{CUMPLE} \\
Consumo humano -- Tabla 1 & $\leq$600 & 0 & \textbf{CUMPLE} \\
\bottomrule
\end{tabular}
\end{table}

Cumple con el límite de TULSMA (CF final = 0 NMP/100mL $≤$ 3000 NMP/100mL).

El producto CT es indicativo de la eficacia de la desinfección. Un valor CT adecuado garantiza la reducción suficiente de microorganismos patógenos.

\begin{table}[H]
\centering
\caption{Verificación del producto CT}
\begin{tabular}{@{}lcc@{}}
\toprule
\textbf{Parámetro} & \textbf{Valor calculado} & \textbf{Valor recomendado} \\
\midrule
CT calculado & 30 mg$\cdot$min/L & $\geq$ 30 mg$\cdot$min/L \\
Estado & \multicolumn{2}{c}{\textbf{Cumple}} \\
\bottomrule
\end{tabular}
\end{table}

El valor CT real (30.0 mg$⋅$min/L) supera el mínimo recomendado de 30 mg$⋅$min/L.


\textbf{Notas de operación}

Se recomienda monitorear el cloro residual en la salida del tanque (debe mantenerse entre 0,5--1,0 mg/L) y ajustar la dosis según la demanda del efluente. Realizar pruebas de coliformes periódicas para verificar la eficacia del sistema.

La reducción logarítmica alcanzada es de 6.6 log, lo que representa una reducción porcentual del 100.0\% de los coliformes presentes en el efluente.

\subsubsection{Resultados}

\begingroup
\small
\begin{longtable}{@{}lc@{}}
\caption{Resumen de dimensiones y parámetros de la desinfección}\\
\toprule
\textbf{Parámetro} & \textbf{Valor} \\
\midrule
\endfirsthead
\caption[]{Resumen de dimensiones y parámetros de la desinfección (continuación)}\\
\toprule
\textbf{Parámetro} & \textbf{Valor} \\
\midrule
\endhead
\midrule
\multicolumn{2}{r}{\textit{Continúa en la siguiente página}} \\
\endfoot
\bottomrule
\endlastfoot
\multicolumn{2}{l}{\textit{Dimensiones del tanque}} \\
\midrule
Largo & 3.0 m \\
Ancho & 3.1 m \\
Canales del culebrin & 5 pasos de 0.50 m \\
Recorrido hidraulico & 15.0 m (30.0:1) \\
Profundidad útil & 2.0 m \\
Altura total & 2.30 m \\
Volumen útil & 18.6 m³ \\
\midrule
\multicolumn{2}{l}{\textit{Parámetros de desinfección}} \\
\midrule
Demanda de cloro & 3.5 mg/L \\
Cloro residual & 0.6 mg/L \\
Dosis total & 4.1 mg/L \\
TRH base & 30 min \\
TRH real & 53.7 min \\
CT & 30 mg$\cdot$min/L \\
Log reducción & 6.6 log \\
CF final & 0 NMP/100mL \\
\midrule
\multicolumn{2}{l}{\textit{Consumo de productos}} \\
\midrule
Consumo Cl$_2$ activo & 2.02 kg/d \\
Concentración NaOCl & 10 \% \\
Consumo NaOCl & 20.2 kg/d \\
Volumen NaOCl & 18.4 L/d (552 L/mes) \\
Almacenamiento (30 d) & 552 L (0.6 m³) \\
\end{longtable}
\endgroup

El sistema de desinfección con hipoclorito de sodio ha sido dimensionado para tratar un caudal de 498.5 m³/d, logrando una reducción de 100.0\% de coliformes fecales. El cumplimiento respecto a los límites de la TULSMA depende de la categoría de uso receptor específica del proyecto, conforme al dictamen técnico presentado en la tabla de verificación anterior.\

\begin{figure}[H]
\centering
\includegraphics[width=0.95\textwidth]{figuras/Esquema_Camara_Contacto_Cloro.png}
\caption{Cámara de contacto de cloración tipo culebrín: planta con bafles alternados y sección transversal. Recorrido hidráulico: 15.0 m; relación recorrido/ancho: 30.0:1; volumen útil adoptado: 18.6 m³.}
\label{fig:camara_contacto_cloro}
\end{figure}



La camara de contacto se configura como culebrin con 5 pasos hidraulicos, ancho de canal de 0.50 m y recorrido total de 15.0 m. La relacion recorrido/ancho es 30.0:1, frente al criterio minimo adoptado de 30:1 para aproximar flujo piston y reducir cortocircuitos hidraulicos. El estado de la verificacion es Cumple.



\subsection{Lechos de Secado de Lodos}



El lecho de secado es una unidad de manejo de lodos que utiliza procesos fisicos de drenaje gravitacional y evaporacion para reducir el contenido de humedad de los lodos generados en el tratamiento. Este sistema es ampliamente utilizado en plantas de tratamiento de pequeno y mediano tamano debido a su bajo consumo energetico y simplicidad operativa.

\textbf{Mecanismos de secado}

El proceso de deshidratación de lodos en el lecho se realiza mediante dos mecanismos fundamentales que actúan simultáneamente: el drenaje gravitacional y la evaporación. El drenaje gravitacional ocurre durante las primeras horas o días después de la aplicación del lodo, cuando el agua libre drena por gravedad a través del medio filtrante (arena y grava) hacia el sistema de drenaje inferior, removiendo aproximadamente del 30 al 50\% del agua inicial del lodo. Una vez agotado el drenaje libre, la humedad restante se pierde por evaporación al ambiente, favorecida por la exposición superficial del lodo y el clima local. En condiciones tropicales de Galápagos, con temperatura promedio de 24°C e insolación elevada, la evaporación es el mecanismo dominante durante la mayor parte del ciclo de secado.

\textit{Nota:} La literatura técnica (Metcalf \& Eddy, 2014; OPS/CEPIS, 2005) reporta proporciones variables entre estos mecanismos según el clima local, con rangos típicos de 1/3 drenaje y 2/3 evaporación en climas cálidos y secos, hasta proporciones inversas en climas húmedos. Para el diseño del presente proyecto se considera el efecto combinado de ambos mecanismos sin asignar proporciones fijas, dado que el tiempo de secado adoptado (15 días) se basa en criterios operativos conservadores de la región.

\textbf{Estructura física del lecho}

El lecho de secado consiste en una estructura impermeabilizada con capas estratificadas que permiten el drenaje y soporte del lodo. La capa superior de arena, con un espesor de 0.25 m, constituye el medio de filtración principal donde se aplica el lodo. La arena fina retiene las partículas de lodo mientras permite el paso del agua, y debe tener granulometría uniforme, típicamente entre 0.25 y 1.0 mm, para maximizar la retención de sólidos sin colmatarse rápidamente. Debajo de esta, la capa intermedia de grava tiene un espesor de 0.20 m y proporciona soporte estructural a la capa de arena, distribuyendo uniformemente el agua drenada hacia el sistema de drenaje. La granulometría de la grava es mayor que la arena, típicamente entre 10 y 30 mm, para crear un medio permeable y estable. Finalmente, en la parte inferior, el sistema de drenaje está compuesto por tuberías perforadas o zanjas que recolectan el agua filtrada y la conducen a un punto de descarga. El gradiente de drenaje debe ser suficiente, con un mínimo del 1 al 2\%, para evitar estancamiento y garantizar el flujo continuo.

\textbf{Ciclo operativo y destino del lodo}

El lecho opera mediante ciclos alternados de carga y descarga. En la fase de aplicación, el lodo húmedo se distribuye uniformemente sobre la superficie del lecho con un espesor de aproximadamente 0.30 m. Durante la fase de secado, el lodo permanece en reposo durante 15 días, período durante el cual pierde humedad por drenaje y evaporación. Finalmente, en la fase de retiro, una vez alcanzado el contenido de humedad objetivo, típicamente entre 40 y 60\% de sólidos, el lodo seco se retira manualmente o con equipos ligeros.

El lodo seco resultante tiene dos destinos principales según su calidad y la normativa local. La disposición en relleno sanitario es la opción más común para lodos de PTAR municipales, ya que el lodo estabilizado y deshidratado cumple con los criterios de manejabilidad para disposición final controlada. Alternativamente, el uso agrícola condicionado es posible solo si el lodo cumple con requisitos estrictos de estabilización, ausencia de metales pesados y patógenos reducidos, según la normativa sanitaria aplicable del Ministerio de Salud y Ministerio del Ambiente, lo cual requiere análisis periódicos y autorización expresa de la autoridad competente.

\textbf{Parámetro de diseño: Carga superficial de sólidos ($\rho_S$)}

El parámetro que controla el área requerida del lecho es la carga superficial de sólidos ($\rho_S$), definida como la masa de sólidos totales (SST) aplicada por unidad de área del lecho por año de operación. Este parámetro integra la producción de lodos con la capacidad de secado del sistema:

\begin{equation}
\rho_S = \frac{M_{SST} \times 365}{A_{total}}
\end{equation}
\captionequation{Carga superficial de solidos -- parametro de control del area}

\textit{Donde:}
\begin{itemize}[noitemsep,leftmargin=2em]
    \item[$\rho_S$] = Carga superficial de sólidos (kg SST/m²·año)
    \item[$M_{SST}$] = Masa diaria de sólidos a tratar (kg SST/d)
    \item[$A_{total}$] = Área superficial total del lecho (m²)
\end{itemize}

La carga superficial $\rho_S$ es crítica porque determina la frecuencia de aplicación de lodos y la eficiencia del secado. Valores elevados de $\rho_S$ implican mayor frecuencia de carga y mayor demanda de área. Los rangos típicos según Metcalf \& Eddy (2014) y OPS/CEPIS (2005) son 60--220 kg SST/m²·año, dependiendo del tipo de lodo, clima y disponibilidad de área.

\textbf{Criterios de diseno}

\begin{table}[H]
\centering
\caption{Parametros de diseno del lecho de secado}
\small
\begin{tabular}{lccl}
\toprule
Parametro & Rango recomendado & Valor adoptado & Fuente \\
\midrule
Carga superficial de solidos & 60--220 kg SST/m²·año & 146 kg/m$^2$·anio & Metcalf \& Eddy \cite{metcalf2014} \\
Concentracion de lodos & 15--30 kg/m³ & 20 kg/m$^3$ & OPS/CEPIS \cite{ops2005} \\
Tiempo de secado & 10--30 días & 15 dias & OPS/CEPIS \\
Espesor de aplicacion & 0.20--0.40 m & 0.30 m & Metcalf \& Eddy \\
Relacion Largo/Ancho & 2:1--4:1 & 3:1 & Practica comun \\
\bottomrule
\end{tabular}
\end{table}

\textbf{Produccion de lodos}

La produccion de lodos se determina del balance encadenado del tren de tratamiento. En este caso, los lodos provienen de: lodos anaerobios del Reactor Anaerobio con Pantallas (ABR/RAP) (factor 0.10 kg SST/kg DBO removida), lodos del Biofiltro de Carga Mecanizada e Hidráulica y lodos del Sedimentador Secundario. Los factores de produccion especificos de cada unidad se aplican segun su tipologia: reactores anaerobios (UASB) tipicamente 0.10 kg SST/kg DBO removida, sistemas aerobios con biopelicula (filtros percoladores, humedales) tipicamente 0.15 kg SST/kg DBO removida, y sedimentadores secundarios segun la produccion de humus estimada.

\textbf{Balance de produccion de lodos:}

\begin{table}[H]
\centering
\small
\caption{Produccion de lodos por linea y total planta}
\begin{tabular}{p{6cm}cc}
\toprule
\textbf{Origen} & \textbf{Por linea} & \textbf{Total planta} \\
 & \textbf{(kg SST/d)} & \textbf{(kg SST/d)} \\
\midrule
Reactor Anaerobio con Pantallas (ABR/RAP) & 0.10 & 0.30 \\
Biofiltro de Carga Mecanizada e Hidráulica & 0.80 & 2.40 \\
Sedimentador Secundario & 0.80 & 2.40 \\
\midrule
\textbf{Total} & \textbf{1.70} & \textbf{5.10} \\
\bottomrule
\end{tabular}
\end{table}

\textbf{Ecuaciones de diseno}

El volumen diario de lodo a tratar se determina mediante:

\begin{equation}
V_{lodo} = \frac{M_{SST}}{C_{SST}}
\end{equation}
\captionequation{Volumen diario de lodo a tratar}

\textit{Donde:}
\begin{itemize}[noitemsep,leftmargin=2em]
    \item[$V_{lodo}$] = Volumen diario de lodo (m$^3$/d)
    \item[$M_{SST}$] = Masa de solidos producidos TOTAL (5.10 kg SST/d)
    \item[$C_{SST}$] = Concentracion de solidos (20 kg/m$^3$)
\end{itemize}

\begin{equation}
V_{lodo} = \frac{5.10}{20} = 0.255 \text{ m}^3\text{/d}
\end{equation}

\textbf{Dimensionamiento del area}

El area requerida se calcula considerando el volumen de lodo, el tiempo de secado y el espesor de aplicacion:

\begin{equation}
A_{total} = \frac{V_{lodo} \cdot t_s}{h_{lodo}}
\end{equation}
\captionequation{Area superficial total del lecho de secado}

\textit{Donde:}
\begin{itemize}[noitemsep,leftmargin=2em]
    \item[$A_{total}$] = Area superficial total requerida (m$^2$)
    \item[$t_s$] = Tiempo de secado (15 dias)
    \item[$h_{lodo}$] = Espesor de aplicacion (0.30 m)
    \item[$n_{celdas}$] = Numero de celdas (3) (subdivision interna)
\end{itemize}

\begin{equation}
A_{total} = \frac{0.255 \times 15}{0.30} = 12.8 \text{ m}^2
\end{equation}

\textbf{Distribucion por lineas de tratamiento}

El area total se distribuye en bloques independientes, uno por cada linea de tratamiento:

\begin{equation}
A_{bloque} = \frac{A_{total}}{n_{lineas}} = \frac{12.8}{3} = 4.3 \text{ m}^2 \text{ por bloque}
\end{equation}

\textbf{Subdivision interna en celdas}

Cada bloque se subdivide internamente en celdas para operacion por rotacion:

\begin{equation}
A_{celda} = \frac{A_{total}}{n_{celdas}} = \frac{12.8}{3} = 4.3 \text{ m}^2 \text{ por celda}
\end{equation}

Con una relacion largo/ancho de 3:1, las dimensiones de cada bloque son 3.6 m de largo por 1.2 m de ancho.

\subsubsection{Verificacion}

\textbf{Verificacion de carga superficial}

La carga superficial de solidos se verifica mediante:

\begin{equation}
\rho_S = \frac{M_{SST} \times 365}{A_{total}} = \frac{5.10 \times 365}{12.8} = 146.0 \text{ kg SST/m}^2\text{·anio}
\end{equation}

El valor obtenido (146.0 kg SST/m²·año) se encuentra dentro del rango recomendado de 60--220 kg SST/m²·año para lechos de secado según Metcalf & Eddy.

\textbf{Notas de operacion}

Se recomienda operar los lechos en sistema de rotación, aplicando lodo en una celda mientras las demás están en proceso de secado o reposo. Bajo operación normal, el lodo secado puede alcanzar contenidos de humedad del 40--60\%, facilitando su manejo y disposición final.

\subsubsection{Resultados}

\begingroup
\small
\begin{longtable}{ll}
\caption{Resumen del dimensionamiento del lecho de secado}\\
\toprule
Parametro & Valor \\
\midrule
\endfirsthead
\caption[]{Resumen del dimensionamiento del lecho de secado (continuación)}\\
\toprule
Parametro & Valor \\
\midrule
\endhead
\midrule
\multicolumn{2}{r}{\textit{Continúa en la siguiente página}} \\
\endfoot
\bottomrule
\endlastfoot
Area total requerida & 12.8 m$^2$ \\
Numero de bloques (corresponde al numero de lineas) & 3 \\
Area por bloque & 4.3 m$^2$ \\
Dimensiones de cada bloque (L$\times$A) & 3.6 m $\times$ 1.2 m \\
Numero total de celdas & 3 \\
Tiempo de secado & 15 dias \\
Carga de solidos & 146.0 kg SST/m$^2$·anio \\
Produccion total de lodos & 5.10 kg SST/d \\
\end{longtable}
\endgroup

\begin{figure}[H]
\centering
\includegraphics[width=0.95\textwidth]{figuras/Esquema_Lecho_Secado.png}
\caption{Lecho de secado de lodos: vista en planta de los bloques por linea de tratamiento y seccion transversal tipica mostrando capas de arena, grava y sistema de drenaje.}
\label{fig:lecho_secado}
\end{figure}

El lecho de secado ha sido dimensionado para tratar 5.10 kg SST/d de lodos generados en el proceso de tratamiento, distribuidos en 3 bloques independientes (uno por linea de tratamiento). Cada bloque opera con un ciclo de 15 dias de secado.

\newpage
\section{Resultados}

\subsection{Disposición de la Planta y Áreas de Predio}

La siguiente figura presenta la disposición espacial de las unidades de tratamiento. El layout muestra 3 líneas paralelas operativas, cada una con capacidad para tratar 5.8~L/s, permitiendo la operación con una sola línea durante mantenimiento o reparaciones. Las dimensiones del área de tratamiento son aproximadamente 51.1~m de ancho por 31.1~m de largo.

\begin{figure}[H]
\centering
\includegraphics[width=\textwidth]{figuras/Layout_sistema_de_tratamiento_abr_rap__taf_ss_3lineas.png}
\caption{Distribución de unidades: \textbf{1}=Rejillas; \textbf{2}=Desarenador; \textbf{3}=Reactor Anaerobio con Pantallas (ABR / RAP); \textbf{4}=Biofiltro de Carga Mecanizada e Hidráulica (TAF); \textbf{5}=Sedimentador Secundario; \textbf{6}=Cámara de Desinfección con Cloro; \textbf{7}=Lecho de Secado de Lodos.}
\label{fig:layout_general}
\end{figure}

El área calculada de 1589~m$^2$ corresponde únicamente a las unidades de tratamiento con sus márgenes mínimos de acceso. Sin embargo, para garantizar una operación segura y eficiente durante toda la vida útil de la planta, es necesario prever espacios adicionales para acceso vehicular y peatonal de operarios y visitantes, almacenamiento seguro de productos químicos y herramientas, realización de análisis de control de calidad in-situ, estacionamiento para personal y visitantes, circulación interna de vehículos de mantenimiento y retiro de lodos, separación de límites con zonas verdes para control de olores, y espacios administrativos para control operativo.

\begin{table}[H]
\centering
\caption{Áreas complementarias de la planta}
\small
\begin{tabular}{lp{8cm}c}
\toprule
Área & Descripción & Dimensión aprox. \\
\midrule
Zona de amortiguación & Perimetral alrededor de unidades (2--3~m) para acceso y seguridad & 318~m$^2$ (20\%) \\
\addlinespace
Bodega de químicos & Almacenamiento de hipoclorito y productos químicos & 19~m$^2$ \\
\addlinespace
Laboratorio/control & Análisis de calidad del agua (DBO, SST, CF, pH) & 29~m$^2$ \\
\addlinespace
Caseta de operación & Control, panel, escritorio, baño & 24~m$^2$ \\
\addlinespace
Área de lavado & Limpieza de equipos y herramientas & 16~m$^2$ \\
\addlinespace
Estacionamiento & Vehículos del personal y visitantes (2--3 vehículos) & 87~m$^2$ \\
\addlinespace
Zona de camiones & Acceso de camiones cisterna/retiro de lodo & 103~m$^2$ \\
\addlinespace
Caminos internas & Circulación vehicular y peatonal (ancho 3--4~m) & 159~m$^2$ (10\% área) \\
\addlinespace
Acceso principal & Entrada, portón y caseta de guardia & 40~m$^2$ \\
\addlinespace
Bodega general & Herramientas, repuestos, EPP & 29~m$^2$ \\
\addlinespace
Zona carga de lodos & Carga/descarga de lodo deshidratado & 40~m$^2$ \\
\addlinespace
Área verde/buffer & Separación de límites, revegetación & 407~m$^2$ (15\% total) \\
\bottomrule
\end{tabular}
\end{table}

\textbf{Área total estimada del predio}

Considerando el área de tratamiento (1589~m$^2$), amortiguación (20\%), complementarias operativas y zona verde (15\% del total):

\begin{equation}
A_{total} = \frac{A_{tratamiento} + A_{amortiguación} + A_{complementarias}}{1 - 0.15}
\end{equation}

\begin{itemize}[noitemsep,leftmargin=2em]
    \item[$A_{tratamiento}$] = 1589~m$^2$
    \item[$A_{amortiguación}$] = 318~m$^2$ (20\%)
    \item[$A_{complementarias}$] = 397~m$^2$ (25\% operativas estimado)
\end{itemize}

\begin{equation}
A_{total} = \frac{1589 + 318 + 397}{0.85} \approx \mathbf{2711 \text{ m}^2} \approx \mathbf{0.27 \text{ ha}}
\end{equation}

\textbf{Nota:} El área total del predio debe ser de aproximadamente 2711~m$^2$ (0.27~ha) para operación adecuada con circulación interna, estacionamiento y zonas verdes.

\newpage



El presente resumen consolida los resultados del dimensionamiento de la Planta de Tratamiento de Aguas Residuales (PTAR) para un caudal de dise\~no de \textbf{17.3 L/s} (5.8 L/s por linea).

\subsection{Parametros de Dise\~no}

\begin{table}[H]
\centering
\caption{Parametros de entrada del sistema}
\begin{tabular}{lc}
\toprule
\textbf{Parametro} & \textbf{Valor} \\
\midrule
Caudal total (Q$_{total}$) & 17.3 L/s (1495.6 m$^3$/d) \\
Caudal por línea (Q$_{linea}$) & 5.8 L/s (498.5 m$^3$/d) \\
DBO$_5$ afluente & 243.1 mg/L \\
DQO afluente & 498.0 mg/L \\
SST afluente & 156.0 mg/L \\
Temperatura & 25.6 °C \\
\bottomrule
\end{tabular}
\end{table}

\subsection{Dimensionamiento de Unidades}

\begin{table}[H]
\centering
\caption{Resumen de unidades de tratamiento}
\small
\begin{tabular}{p{3.2cm}p{3.5cm}p{1.5cm}p{4.8cm}}
\toprule
\textbf{Unidad} & \textbf{Dimensiones (m)} & \textbf{Cantidad} & \textbf{Parametro Clave} \\
\midrule
Rejillas & 0.60 $\times$ 0.40 & 3 & Velocidad: 0.50 m/s \\
Desarenador & 0.60 $\times$ 3.0 $\times$ 0.7 & 3 & TRH: 125~s (referencia: 30--60~s) \\
ABR-RAP & 10.7 $\times$ 1.79 & 3 & v$_{up}$ = 1.94 m/h \\
TAF & D = 3.50, H = 4.0 & 3 & CHS = 4.33 m$^3$/m$^2$·h \\
Sedimentador Secundario & D = 6.30, H = 3.5 & 3 & SOR = 16.0 m$^3$/m$^2\cdot$d \\
Desinfeccion (Cloro) & 3.0 $\times$ 3.1 $\times$ 2.3 & 3 & CT = 30 mg$\cdot$min/L \\
Lecho de Secado & 3.6 $\times$ 1.2 & 3 & Área: 4.3 m$^2$ \\
\bottomrule
\end{tabular}
\end{table}

\subsection{Balance de Calidad del Agua}

\begin{table}[H]
\centering
\caption{Evolucion de parametros de calidad a traves del tratamiento}
\begin{tabular}{lccccc}
\toprule
\textbf{} & \textbf{Afluente} & \textbf{Post-ABR} & \textbf{Post-TAF} & \textbf{Post-Sed} & \textbf{Efluente} \\
\midrule
DBO$_5$ (mg/L) & 243.1 & 57.7 (75\%) & 25.6 (56\%) & 21.8 (15\%) & 21.8 \\
DQO (mg/L) & 498.0 & 141.9 (70\%) & 70.9 (50\%) & 60.2 (15\%) & 60.2 \\
SST (mg/L) & 156.0 & 31.5 (75\%) & 48.5 (0\%) & 9.7 (80\%) & 9.7 \\
CF (NMP/100mL) & 3300 & 1787 (40\%) & 1430 (20\%) & 1287 (10\%) & 1 \\
\bottomrule
\end{tabular}
\end{table}

\subsection{Verificacion de Cumplimiento TULSMA}

\textbf{Efluente calculado:} DBO$_5$ = 21.8 mg/L | DQO = 60.2 mg/L | SST = 9.7 mg/L | CF = 1 NMP/100mL

\vspace{0.5em}
\begin{table}[H]
\centering
\small
\caption{Comparacion con limites TULSMA por uso del agua}
\begin{tabular}{lcccccc}
\toprule
\textbf{Uso / Tabla TULSMA} & \textbf{DBO$_5$} & \textbf{DQO} & \textbf{SST} & \textbf{CF} & \textbf{Dictamen} \\
& \textbf{(limite)} & \textbf{(limite)} & \textbf{(limite)} & \textbf{(limite)} & \\
\midrule
Agua dulce -- Tabla 12 & $\leq$100 & $\leq$250 & $\leq$100 & $\leq$3000 & CUMPLE \\
Agua marina -- Tabla 13 & $\leq$100 & $\leq$250 & $\leq$100 & $\leq$3000 & CUMPLE \\
Alcantarillado -- Tabla 11 & $\leq$250 & $\leq$500 & $\leq$220 & --- & CUMPLE \\
Flora y fauna -- Tabla 3 & --- & --- & --- & $\leq$200 & CUMPLE \\
Agricola/riego -- Tabla 6 & --- & --- & --- & $\leq$1000 & CUMPLE \\
Pecuario -- Tabla 7 & --- & --- & --- & $\leq$1000 & CUMPLE \\
Recreativo primario -- Tabla 9 & --- & --- & --- & $\leq$200 & CUMPLE \\
Recreativo secundario -- Tabla 10 & --- & --- & --- & $\leq$2000 & CUMPLE \\
Consumo humano trat. conv. -- Tabla 1 & $\leq$2.0 & --- & --- & $\leq$600 & NO CUMPLE \\
\bottomrule
\end{tabular}
\end{table}

\subsection{Requerimientos de Terreno}

\begin{table}[H]
\centering
\caption{Resumen de areas requeridas}
\begin{tabular}{lc}
\toprule
\textbf{Concepto} & \textbf{Area} \\
\midrule
Area de tratamiento & 1589 m$^2$ \\
Area de amortiguacion (20\%) & 318 m$^2$ \\
Area complementaria operativa & 397 m$^2$ (25\% estimado) \\
Zona verde (15\% del total) & 407 m$^2$ \\
\midrule
\textbf{Area total requerida} & \textbf{2711 m$^2$ (0.27 ha)} \\
\bottomrule
\end{tabular}
\end{table}

\subsection{Produccion de Lodos}

\begin{table}[H]
\centering
\caption{Generacion y manejo de lodos - TOTAL PLANTA (3 líneas)}
\begin{tabular}{p{6cm}cc}
\toprule
\textbf{Concepto} & \textbf{Por línea} & \textbf{Total planta} \\
\midrule
Reactor Anaerobio con Pantallas (ABR/RAP) & 0.10 kg SST/d & 0.30 kg SST/d \\
Biofiltro de Carga Mecanizada e Hidráulica & 0.80 kg SST/d & 2.40 kg SST/d \\
Sedimentador Secundario & 0.80 kg SST/d & 2.40 kg SST/d \\
\midrule
\textbf{Total lodos} & \textbf{1.70 kg SST/d} & \textbf{5.10 kg SST/d} \\
Area de lechos de secado & --- & 12.8 m$^2$ (3 bloques) \\
Frecuencia de aplicacion & --- & 1 vez cada 3--4 meses \\
\bottomrule
\end{tabular}
\end{table}

\subsection{Consumos Estimados}

\begin{table}[H]
\centering
\caption{Consumos de energia y productos quimicos}
\begin{tabular}{lcc}
\toprule
\textbf{Concepto} & \textbf{Consumo} & \textbf{Unidad} \\
\midrule
Potencia electrica estimada & 5--8 & kW \\
Consumo energia & 40,000--60,000 & kWh/a\~no \\
Hipoclorito de sodio (10\%) & 18.4 & L/d \\
Consumo hipoclorito anual & 7 & m$^3$/a\~no \\
\bottomrule
\end{tabular}
\end{table}




\newpage
\bibliographystyle{plain}
\bibliography{referencias}


\newpage
\appendix
\section{Anexo: TULSMA -- Criterios de Calidad del Recurso Agua}
\label{app:tulsma}

\begin{center}
{\LARGE\bfseries TULSMA -- Norma de Calidad Ambiental\\[4pt]
del Recurso Agua}\\[8pt]
{\large Libro VI, Anexo 1 -- Acuerdo Ministerial 097-A (2015)\\
Decreto Ejecutivo 3516 -- Ecuador}\\[4pt]
{\small Criterios de calidad por uso del agua y limites de descarga de efluentes}
\end{center}
\vspace{10pt}

\subsection{Tabla 1 -- Consumo Humano y Uso Domestico (tratamiento convencional)}
\small
\begin{longtable}{p{4.5cm}p{2.0cm}p{2.0cm}p{3.5cm}}
\caption{Limites maximos permisibles para aguas de consumo humano y uso domestico que unicamente requieren tratamiento convencional (TULSMA, Tabla~1).}\\
\toprule
\textbf{Parametro} & \textbf{Expresado como} & \textbf{Unidad} & \textbf{Limite maximo permisible} \\
\midrule
\endfirsthead
\multicolumn{4}{c}{\small\itshape Continuacion de la Tabla}\\
\toprule
\textbf{Parametro} & \textbf{Expresado como} & \textbf{Unidad} & \textbf{Limite maximo permisible} \\
\midrule
\endhead
\bottomrule
\endlastfoot
Aceites y Grasas & Sust.\ solubles en hexano & mg/L & $\leq$ 0,3 \\
Aluminio & Al & mg/L & $\leq$ 0,2 \\
Amoniaco & N-Amoniacal & mg/L & $\leq$ 1,0 \\
Amonio & NH$_4$ & mg/L & $\leq$ 0,05 \\
Arsenico (total) & As & mg/L & $\leq$ 0,05 \\
Bario & Ba & mg/L & $\leq$ 1,0 \\
Cadmio & Cd & mg/L & $\leq$ 0,01 \\
Cianuro total & CN$^-$ & mg/L & $\leq$ 0,1 \\
Cloruros & Cl$^-$ & mg/L & $\leq$ 250 \\
Cobre & Cu & mg/L & $\leq$ 1,0 \\
\textbf{Coliformes totales} & NMP & NMP/100 mL & $\leq$ 3\,000 \\
\textbf{Coliformes fecales} & NMP & NMP/100 mL & $\leq$ 600 \\
Color real & UC & Unidades de color & $\leq$ 100 \\
Compuestos fenolicos & Fenol & mg/L & $\leq$ 0,002 \\
Cromo hexavalente & Cr$^{6+}$ & mg/L & $\leq$ 0,05 \\
\textbf{DBO$_5$} & -- & mg/L & $\leq$ 2,0 \\
Dureza & CaCO$_3$ & mg/L & $\leq$ 500 \\
Fluoruro total & F & mg/L & $\leq$ 1,5 \\
Hierro total & Fe & mg/L & $\leq$ 1,0 \\
Manganeso total & Mn & mg/L & $\leq$ 0,1 \\
Materia flotante & -- & -- & Ausencia \\
Mercurio total & Hg & mg/L & $\leq$ 0,001 \\
Nitrato & N-Nitrato & mg/L & $\leq$ 10,0 \\
Nitrito & N-Nitrito & mg/L & $\leq$ 1,0 \\
Oxigeno disuelto & O.D. & mg/L & $\geq$ 80\% saturacion; $\geq$ 6 mg/L \\
pH & -- & -- & $6,0 - 9,0$ \\
Plata total & Ag & mg/L & $\leq$ 0,05 \\
Plomo total & Pb & mg/L & $\leq$ 0,05 \\
Selenio total & Se & mg/L & $\leq$ 0,01 \\
Sodio & Na & mg/L & $\leq$ 200 \\
Solidos disueltos totales & -- & mg/L & $\leq$ 1\,000 \\
Sulfatos & SO$_4^{2-}$ & mg/L & $\leq$ 400 \\
Temperatura & -- & $^\circ$C & Cond.\ natural $\pm$ 3\,$^\circ$C \\
Tensoactivos & SAAM & mg/L & $\leq$ 0,5 \\
Turbiedad & -- & UTN & $\leq$ 100 \\
Zinc & Zn & mg/L & $\leq$ 5,0 \\
\end{longtable}

\subsection{Tabla 12 -- Descarga a Cuerpo de Agua Dulce}
{\small\textbf{Nota:} Esta es la tabla aplicable al proyecto de tratamiento de aguas residuales municipales cuando el efluente se vierte a un rio, lago o cuerpo superficial de agua dulce.}
\vspace{6pt}
\small
\begin{longtable}{p{4.5cm}p{2.0cm}p{2.0cm}p{3.5cm}}
\caption{Limites maximos permisibles de descarga a un cuerpo de agua dulce (TULSMA, Tabla~12).}\\
\toprule
\textbf{Parametro} & \textbf{Expresado como} & \textbf{Unidad} & \textbf{Limite maximo permisible} \\
\midrule
\endfirsthead
\multicolumn{4}{c}{\small\itshape Continuacion de la Tabla}\\
\toprule
\textbf{Parametro} & \textbf{Expresado como} & \textbf{Unidad} & \textbf{Limite maximo permisible} \\
\midrule
\endhead
\bottomrule
\endlastfoot
Aceites y grasas & Sust.\ solubles en hexano & mg/L & $\leq$ 0,3 \\
Aluminio & Al & mg/L & $\leq$ 5,0 \\
Arsenico total & As & mg/L & $\leq$ 0,1 \\
Bario & Ba & mg/L & $\leq$ 2,0 \\
Cadmio & Cd & mg/L & $\leq$ 0,02 \\
Cianuro total & CN$^-$ & mg/L & $\leq$ 0,1 \\
Cloruros & Cl$^-$ & mg/L & $\leq$ 1\,000 \\
Cobre & Cu & mg/L & $\leq$ 1,0 \\
\textbf{Coliformes fecales}$^\dagger$ & NMP & NMP/100 mL & $\leq$ 3\,000 \\
Compuestos fenolicos & Fenol & mg/L & $\leq$ 0,2 \\
Cromo hexavalente & Cr$^{6+}$ & mg/L & $\leq$ 0,5 \\
\textbf{DBO$_5$} & -- & mg/L & $\leq$ 100 \\
\textbf{DQO} & -- & mg/L & $\leq$ 250 \\
Fluoruros & F$^-$ & mg/L & $\leq$ 5,0 \\
Fosforo total & P & mg/L & $\leq$ 10,0 \\
Hierro total & Fe & mg/L & $\leq$ 10,0 \\
Manganeso total & Mn & mg/L & $\leq$ 2,0 \\
Materia flotante & -- & -- & Ausencia \\
Mercurio total & Hg & mg/L & $\leq$ 0,005 \\
Niquel & Ni & mg/L & $\leq$ 2,0 \\
Nitrogeno amoniacal & N-NH$_3$ & mg/L & $\leq$ 30,0 \\
Nitrogeno total Kjeldahl & N-TKN & mg/L & $\leq$ 50,0 \\
pH & -- & -- & $6,0 - 9,0$ \\
Plata total & Ag & mg/L & $\leq$ 0,1 \\
Plomo total & Pb & mg/L & $\leq$ 0,2 \\
Selenio & Se & mg/L & $\leq$ 0,1 \\
\textbf{Solidos sedimentables} & -- & mL/L & $\leq$ 1,0 \\
\textbf{Solidos suspendidos totales} & -- & mg/L & $\leq$ 130 \\
Solidos totales & -- & mg/L & $\leq$ 1\,600 \\
Sulfatos & SO$_4^{2-}$ & mg/L & $\leq$ 1\,000 \\
Sulfuros & S$^{2-}$ & mg/L & $\leq$ 0,5 \\
Temperatura & -- & $^\circ$C & Cond.\ natural $\pm$ 3\,$^\circ$C; max.\ 32\,$^\circ$C \\
Tensoactivos & SAAM & mg/L & $\leq$ 0,5 \\
Zinc & Zn & mg/L & $\leq$ 5,0 \\
\multicolumn{4}{l}{\scriptsize $^\dagger$ Aquellos con descargas $\leq$ 3\,000 NMP/100 mL quedan exentos de tratamiento de desinfeccion.}\\
\end{longtable}
\normalsize

\subsection{Tabla 13 -- Descarga a Cuerpo de Agua Marina}
\small
\begin{longtable}{p{4.5cm}p{2.0cm}p{2.0cm}p{3.5cm}}
\caption{Limites maximos permisibles de descarga a un cuerpo de agua marina (TULSMA, Tabla~13).}\\
\toprule
\textbf{Parametro} & \textbf{Expresado como} & \textbf{Unidad} & \textbf{Limite maximo permisible} \\
\midrule
\endfirsthead
\multicolumn{4}{c}{\small\itshape Continuacion de la Tabla}\\
\toprule
\textbf{Parametro} & \textbf{Expresado como} & \textbf{Unidad} & \textbf{Limite maximo permisible} \\
\midrule
\endhead
\bottomrule
\endlastfoot
Aceites y grasas & Sust.\ solubles en hexano & mg/L & $\leq$ 0,3 \\
Aluminio & Al & mg/L & $\leq$ 5,0 \\
Arsenico total & As & mg/L & $\leq$ 0,1 \\
Cadmio & Cd & mg/L & $\leq$ 0,1 \\
Cianuro total & CN$^-$ & mg/L & $\leq$ 0,1 \\
Cobre & Cu & mg/L & $\leq$ 0,5 \\
\textbf{Coliformes fecales}$^\ddagger$ & NMP & NMP/100 mL & $\leq$ 3\,000 \\
Compuestos fenolicos & Fenol & mg/L & $\leq$ 0,2 \\
Cromo hexavalente & Cr$^{6+}$ & mg/L & $\leq$ 0,5 \\
\textbf{DBO$_5$} & -- & mg/L & $\leq$ 100 \\
\textbf{DQO} & -- & mg/L & $\leq$ 250 \\
Fluoruros & F$^-$ & mg/L & $\leq$ 5,0 \\
Hierro total & Fe & mg/L & $\leq$ 10,0 \\
Manganeso total & Mn & mg/L & $\leq$ 2,0 \\
Materia flotante & -- & -- & Ausencia \\
Mercurio total & Hg & mg/L & $\leq$ 0,005 \\
Niquel & Ni & mg/L & $\leq$ 2,0 \\
Nitrogeno amoniacal & N-NH$_3$ & mg/L & $\leq$ 30,0 \\
Nitrogeno total Kjeldahl & N-TKN & mg/L & $\leq$ 50,0 \\
pH & -- & -- & $6,0 - 9,0$ \\
Plomo total & Pb & mg/L & $\leq$ 0,2 \\
\textbf{Solidos suspendidos totales} & -- & mg/L & $\leq$ 100 \\
Temperatura & -- & $^\circ$C & Cond.\ natural $\pm$ 3\,$^\circ$C; max.\ 32\,$^\circ$C \\
Tensoactivos & SAAM & mg/L & $\leq$ 0,5 \\
Zinc & Zn & mg/L & $\leq$ 5,0 \\
\multicolumn{4}{l}{\footnotesize $^\ddagger$ Aquellos con descargas $\leq$ 3\,000 NMP/100 mL quedan exentos de tratamiento de desinfeccion.}\\
\end{longtable}

\subsection{Tabla 11 -- Descarga al Sistema de Alcantarillado Publico}
\small
\begin{longtable}{p{4.5cm}p{2.0cm}p{2.0cm}p{3.5cm}}
\caption{Limites maximos permisibles de descarga al sistema de alcantarillado publico (TULSMA, Tabla~11).}\\
\toprule
\textbf{Parametro} & \textbf{Expresado como} & \textbf{Unidad} & \textbf{Limite maximo permisible} \\
\midrule
\endfirsthead
\multicolumn{4}{c}{\small\itshape Continuacion de la Tabla}\\
\toprule
\textbf{Parametro} & \textbf{Expresado como} & \textbf{Unidad} & \textbf{Limite maximo permisible} \\
\midrule
\endhead
\bottomrule
\endlastfoot
Aceites y grasas & Sust.\ solubles en hexano & mg/L & $\leq$ 70,0 \\
Aluminio & Al & mg/L & $\leq$ 5,0 \\
Arsenico total & As & mg/L & $\leq$ 0,1 \\
Bario & Ba & mg/L & $\leq$ 2,0 \\
Cadmio & Cd & mg/L & $\leq$ 0,02 \\
Cianuro total & CN$^-$ & mg/L & $\leq$ 1,0 \\
Cloruros & Cl$^-$ & mg/L & $\leq$ 1\,000 \\
Cobre & Cu & mg/L & $\leq$ 1,0 \\
Coliformes fecales & NMP & NMP/100 mL & No especificado \\
Compuestos fenolicos & Fenol & mg/L & $\leq$ 0,2 \\
Cromo hexavalente & Cr$^{6+}$ & mg/L & $\leq$ 0,5 \\
Cromo total & Cr & mg/L & $\leq$ 0,5 \\
\textbf{DBO$_5$} & -- & mg/L & $\leq$ 250 \\
\textbf{DQO} & -- & mg/L & $\leq$ 500 \\
Fluoruros & F$^-$ & mg/L & $\leq$ 5,0 \\
Hierro total & Fe & mg/L & $\leq$ 25,0 \\
Manganeso total & Mn & mg/L & $\leq$ 10,0 \\
Materia flotante & -- & -- & Ausencia \\
Mercurio total & Hg & mg/L & $\leq$ 0,01 \\
Niquel & Ni & mg/L & $\leq$ 2,0 \\
Nitrogeno amoniacal & N-NH$_3$ & mg/L & $\leq$ 30,0 \\
pH & -- & -- & $6,0 - 9,0$ \\
Plata total & Ag & mg/L & $\leq$ 0,5 \\
Plomo total & Pb & mg/L & $\leq$ 0,5 \\
Selenio & Se & mg/L & $\leq$ 0,5 \\
\textbf{Solidos sedimentables} & -- & mL/L/h & $\leq$ 20 \\
\textbf{Solidos suspendidos totales} & -- & mg/L & $\leq$ 220 \\
Sulfatos & SO$_4^{2-}$ & mg/L & $\leq$ 1\,000 \\
Sulfuros & S$^{2-}$ & mg/L & $\leq$ 1,0 \\
Temperatura & -- & $^\circ$C & $\leq$ 40 \\
Tensoactivos & SAAM & mg/L & $\leq$ 2,0 \\
Zinc & Zn & mg/L & $\leq$ 5,0 \\
\end{longtable}

\end{document}